% Luc Destombe --- licence CC BY-NC-SA 4.0
% https://creativecommons.org/licenses/by-nc-sa/4.0/deed.fr
% Fichier autonome : compiler avec pdflatex, deux fois (signets, nombre de pages).
\documentclass[11pt,a4paper]{article}
\makeatletter
\newif\iflycee@unecolonne
\lycee@unecolonnetrue
\newif\iflycee@palatino
\lycee@palatinotrue

\RequirePackage[utf8]{inputenc}
\RequirePackage[T1]{fontenc}
\RequirePackage[french]{babel}
\RequirePackage{etoolbox}

\newcommand{\ShorthandsOff}{\@ifundefined{shorthandoff}{}{\shorthandoff{;:!?}}}
\newcommand{\ShorthandsOn}{\@ifundefined{shorthandon}{}{\shorthandon{;:!?}}}
\ShorthandsOff
\AtBeginDocument{\ShorthandsOn}
\AtBeginEnvironment{tikzpicture}{\ShorthandsOff}

\RequirePackage{geometry}
\RequirePackage{fancyhdr}
\RequirePackage{lastpage}
\RequirePackage{multicol}
\RequirePackage{array}
\RequirePackage{enumitem}
\RequirePackage{xcolor}
\RequirePackage{needspace}

\RequirePackage{amsmath,amssymb}
\RequirePackage{mathpazo}
\DeclareMathSizes{10.95}{10.95}{8}{5.5}
\begingroup\catcode`\_=\active \catcode`\^=\active
\gdef_#1{\sb{\scriptscriptstyle #1}}
\gdef^#1{\sp{\scriptscriptstyle\lycee@moinsepais #1}}
\endgroup
\newcommand\lycee@moins{\mathbin{%
  \setbox\z@\hbox{$\m@th\scriptscriptstyle\lycee@moinsorig$}%
  \dimen@=.55\wd\z@ \advance\dimen@-.5pt
  \hbox to.75\wd\z@{\hss$\m@th\scriptscriptstyle\vcenter{\hbox{%
    \tikz\draw[line width=.5pt,line cap=round](0,0)--(\the\dimen@,0);}}$\hss}}}
\begingroup\lccode`\~=`\- \lowercase{\endgroup\let~}\lycee@moins
\newcommand\lycee@moinsepais{\mathcode`\-="8000 }
\AtBeginDocument{\mathchardef\lycee@moinsorig=\mathcode`\-\relax}
\AtBeginDocument{\catcode`\_=12 \mathcode`\_="8000
  \catcode`\^=12 \mathcode`\^="8000 }
\newcommand\lycee@exposants{%
  \fontdimen13\textfont2=.48\fontdimen6\textfont2
  \fontdimen14\textfont2=.48\fontdimen6\textfont2
  \fontdimen15\textfont2=.48\fontdimen6\textfont2
  \fontdimen13\scriptfont2=.48\fontdimen6\scriptfont2
  \fontdimen14\scriptfont2=.48\fontdimen6\scriptfont2
  \fontdimen15\scriptfont2=.48\fontdimen6\scriptfont2}
\every@math@size\expandafter{\the\every@math@size\lycee@exposants}
\newlength{\retraitcalculs}
\setlength{\retraitcalculs}{2em}

\RequirePackage{tikz}
\usetikzlibrary{arrows.meta,calc,babel,shapes.misc}
\RequirePackage[most]{tcolorbox}
\tcbuselibrary{skins,breakable}

\definecolor{coulDef}{HTML}{1F4E79}
\definecolor{coulProp}{HTML}{7030A0}
\definecolor{coulRem}{HTML}{C55A11}
\definecolor{coulEx}{HTML}{2E7D32}
\definecolor{coulExo}{HTML}{B03A2E}
\definecolor{coulPart}{HTML}{1F4E79}
\definecolor{coulMeth}{HTML}{1B6E4F}
\definecolor{coulCourbe}{HTML}{0B5ED7}
\definecolor{coulCourbeB}{HTML}{C00000}
\definecolor{coulCourbeC}{HTML}{2E7D32}
\definecolor{coulGrille}{HTML}{8C8C8C}
\definecolor{coulRep}{HTML}{1B6E4F}
\definecolor{coulBar}{HTML}{2E7D32}

\newcommand{\R}{\mathbb{R}}
\newcommand{\Rs}{\mathbb{R}^{*}}
\renewcommand{\le}{\leqslant}
\renewcommand{\ge}{\geqslant}
\renewcommand{\approx}{\simeq}
\newcommand{\vect}[1]{\overrightarrow{#1}}
\newcommand{\norme}[1]{\left\lVert #1 \right\rVert}
\newcommand{\intervalle}[2]{\left[#1\,;#2\right]}
\RequirePackage{cancel}
\renewcommand{\CancelColor}{\color{coulExo}}
\newcommand{\fois}[1]{\times{\color{coulExo}#1}}
\newcommand{\barre}[1]{\cancel{#1}}

\tikzset{grille/.style={coulGrille,line width=0.4pt}}
\newcommand{\quadrillage}[4]{\draw[grille,step=1] (#1,#2) grid (#3,#4);}
\tikzset{
  croix/.style={cross out,draw,line width=0.6pt,inner sep=0pt,minimum size=4pt},
  croixdroite/.style={croix,rotate=45,line width=0.8pt},
}
\newcommand{\pt}[3]{\path (#1) node[croix]{} node[#2,font=\small,inner sep=2.5pt]{$#3$};}

\newcommand{\lycee@repere}[4]{%
  \draw[grille,step=1] (#1,#2) grid (#3,#4);
  \draw[-{Stealth[length=2mm]},thick] (#1,0)--({#3+0.4},0) node[below left=-1pt]{$x$};
  \draw[-{Stealth[length=2mm]},thick] (0,#2)--(0,{#4+0.4}) node[below left=-1pt]{$y$};
  \node[below left,font=\scriptsize,inner sep=1.5pt] at (0,0) {$0$};}
\newcommand{\lycee@gradx}[1]{\foreach \k/\l in {#1}{%
  \draw (\k,0.07)--(\k,-0.07) node[below,font=\scriptsize,inner sep=1.5pt]{$\l$};}}
\newcommand{\lycee@grady}[1]{\foreach \k/\l in {#1}{%
  \draw (0.07,\k)--(-0.07,\k) node[left,font=\scriptsize,inner sep=1.5pt]{$\l$};}}
\newcommand{\lycee@intff}[2]{\left[#1\,;#2\right]}
\newcommand{\lycee@intfo}[2]{\left[#1\,;#2\right[}
\newcommand{\lycee@intof}[2]{\left]#1\,;#2\right]}
\newcommand{\lycee@intoo}[2]{\left]#1\,;#2\right[}
\newcommand{\lycee@ens}[1]{\left\{#1\right\}}
\AtBeginDocument{%
  \providecommand{\repere}{\lycee@repere}%
  \providecommand{\gradx}{\lycee@gradx}%
  \providecommand{\grady}{\lycee@grady}%
  \providecommand{\Cf}{\mathcal{C}\sb{\scriptscriptstyle f}}%
  \providecommand{\Cg}{\mathcal{C}\sb{\scriptscriptstyle g}}%
  \providecommand{\intff}{\lycee@intff}%
  \providecommand{\intfo}{\lycee@intfo}%
  \providecommand{\intof}{\lycee@intof}%
  \providecommand{\intoo}{\lycee@intoo}%
  \providecommand{\ens}{\lycee@ens}}

\newlist{questions}{enumerate}{3}
\setlist[questions,1]{label=\arabic*\textdegree),leftmargin=*,itemsep=2pt,topsep=3pt}
\setlist[questions,2]{label=\alph*),leftmargin=*,itemsep=1pt,topsep=2pt}
\setlist[questions,3]{label=\roman*),leftmargin=*,itemsep=1pt,topsep=1pt}
\setlist[itemize]{label=\textbullet,leftmargin=*,itemsep=2pt,topsep=3pt}

\newcommand{\besoinplace}[1]{\par\penalty\@M\begingroup
  \dimen@=#1\relax\dimen@ii=\pagegoal\advance\dimen@ii-\pagetotal
  \ifdim\dimen@>\dimen@ii\ifdim\dimen@ii>\z@\vfil\fi\break\fi\endgroup}

\newcommand{\lycee@pied}{\thepage/\pageref{LastPage}}
\newcommand{\entete}[2]{%
  \pagestyle{fancy}\fancyhf{}%
  \fancyhead[L]{\footnotesize\color{coulPart}\textbf{#1}}%
  \fancyhead[R]{\footnotesize\color{coulPart}\textbf{#2}}%
  \fancyfoot[C]{\footnotesize\lycee@pied}%
  \renewcommand{\headrulewidth}{0.4pt}%
  \renewcommand{\headrule}{\hbox to\headwidth{%
    \color{coulPart!40}\leaders\hrule height \headrulewidth\hfill}}}

\RequirePackage{ccicons}
\newcommand{\lycee@licence}{{\scriptsize\color{gray}%
  \copyright\ Luc Destombe --- \ccbyncsa\ CC BY-NC-SA 4.0}}
\newif\iflycee@licence \lycee@licencetrue
\newcommand{\sanslicence}{\lycee@licencefalse}
\AtBeginDocument{\iflycee@licence\AddToHookNext{shipout/foreground}{%
  \put(0,0){\rlap{\kern\dimexpr1in+\hoffset+\oddsidemargin+\textwidth\relax
    \llap{\raisebox{-\dimexpr1in+\voffset+\topmargin+\headheight+\headsep
      +\textheight+\footskip\relax}{\lycee@licence}}}}}\fi}

\newcommand{\formulesserrees}{%
  \setlength{\abovedisplayskip}{3pt}\setlength{\belowdisplayskip}{3pt}%
  \setlength{\abovedisplayshortskip}{2pt}\setlength{\belowdisplayshortskip}{3pt}}

\setlength{\parindent}{0pt}

\RequirePackage[implicit=false,hidelinks,pdfencoding=auto,psdextra,bookmarksopen,
  bookmarksopenlevel=1,pdfcreator={},pdfproducer={}]{hyperref}
\RequirePackage{bookmark}
\hypersetup{pdfauthor={Luc Destombe},pdfkeywords={CC BY-NC-SA 4.0}}
\newcounter{lycee@signet}
\newcommand{\signet}[3][]{\stepcounter{lycee@signet}%
  \edef\lycee@dest{\if\relax\detokenize{#1}\relax
    signet.\arabic{lycee@signet}\else#1\fi}%
  \hypertarget{\lycee@dest}{}\bookmark[level=#2,dest=\lycee@dest]{#3}}
\pdfstringdefDisableCommands{%
  \def\R{R}\def\vect#1{#1}\def\Cf{Cf}\def\Cg{Cg}\def\fois#1{×#1}%
  \def\barre#1{#1}\def\intervalle#1#2{[#1;#2]}\def\intff#1#2{[#1;#2]}%
  \def\textsuperscript#1{#1}\def\dfrac#1#2{#1/#2}\def\frac#1#2{#1/#2}%
  \def\vec#1{#1⃗}}%
\RequirePackage[official]{eurosym}

\tcbset{exercice corrige/.style={
  enhanced,breakable,before upper=\raggedright,
  colback=white,colframe=coulExo,
  boxrule=0.8pt,arc=2pt,
  left=6pt,right=6pt,top=8pt,bottom=5pt,
  before skip=9pt,after skip=4pt,
  coltitle=white,fonttitle=\bfseries\small,
  attach boxed title to top left={xshift=8pt,yshift=-\tcboxedtitleheight/2},
  boxed title style={colback=coulExo,arc=2pt,boxrule=0pt}}}
\newcounter{exo}
\newtcolorbox{exercice}[1][]{exercice corrige,
  phantom={\signet[exercice-\arabic{exo}]{2}{Exercice \arabic{exo}}},
  title={Exercice \theexo\ifx\relax#1\relax\else{} --- #1\fi}}
\newenvironment{exo}[1][\relax]{\refstepcounter{exo}\begin{exercice}[#1]}{\end{exercice}}

\geometry{top=1.8cm,bottom=1cm,left=1cm,right=1cm,headheight=14pt,footskip=0.6cm}
\renewcommand{\lycee@pied}{\thepage}

\renewcommand{\theexo}{\ifnum\value{exo}<10 0\fi\arabic{exo}}
\tcbset{exercice corrige/.style={
  enhanced,breakable,before upper=\raggedright,
  colback=white,colframe=coulExo,
  boxrule=0.9pt,arc=2pt,
  left=8pt,right=8pt,top=8pt,bottom=6pt,
  before skip=8pt,after skip=6pt,
  coltitle=white,fonttitle=\bfseries,
  attach boxed title to top left={xshift=8pt,yshift=-\tcboxedtitleheight/2},
  boxed title style={colback=coulExo,arc=2pt,boxrule=0pt}}}

\newcommand{\entetefiche}[2]{%
  \begin{center}
    \begin{tcolorbox}[enhanced,width=0.92\linewidth,colback=white,colframe=coulPart,
        boxrule=1.4pt,arc=3pt,halign=center,top=6pt,bottom=6pt]
      {\large\bfseries\color{coulPart} CORRIGÉ --- FICHE D'EXERCICES}\\[3pt]
      {\normalsize\bfseries #1}\\[2pt]
      {\normalsize #2}
    \end{tcolorbox}
  \end{center}}

\newcounter{partie}
\newcommand{\partie}[1]{%
  \stepcounter{partie}%
  \besoinplace{8\baselineskip}\vspace{4pt}%
  \begin{tcolorbox}[enhanced,colback=coulPart!8,colframe=coulPart,boxrule=0pt,
      leftrule=3pt,arc=0pt,left=6pt,right=4pt,top=3pt,bottom=3pt,
      before skip=6pt,after skip=2pt]
    \bfseries\color{coulPart}\Roman{partie} --- #1\signet{1}{\Roman{partie} --- #1}
  \end{tcolorbox}}
\newcommand{\partiesansnum}[1]{%
  \besoinplace{8\baselineskip}\vspace{4pt}%
  \begin{tcolorbox}[enhanced,colback=coulPart!8,colframe=coulPart,boxrule=0pt,
      leftrule=3pt,arc=0pt,left=6pt,right=4pt,top=3pt,bottom=3pt,
      before skip=6pt,after skip=2pt]
    \bfseries\color{coulPart}#1\signet{1}{#1}
  \end{tcolorbox}}

\setlist[questions,1]{label=\arabic*\textdegree),leftmargin=*,itemsep=3pt,topsep=3pt}
\setlist[questions,2]{label=\alph*),leftmargin=*,itemsep=2pt,topsep=2pt}
\setlist[questions,3]{label=\roman*),leftmargin=*,itemsep=1pt,topsep=1pt}
\newlist{expr}{enumerate}{1}
\setlist[expr]{label={\Alph*\ $=$},leftmargin=*,itemsep=2pt,topsep=3pt}

\newcommand{\res}[1]{{\color{coulRep}#1}}
\newcommand{\rem}[1]{\par\smallskip{\small\itshape\color{black!60}$\rightarrow$ #1}}
\newenvironment{calculs}{\par\smallskip\noindent$\displaystyle\begin{aligned}[t]}%
  {\end{aligned}$\par\smallskip}
\newcommand{\justif}[1]{\quad\text{\small\itshape\color{black!60}#1}}
\newcommand{\cas}[1]{\par\smallskip\noindent\textbf{\color{coulExo}#1}\nobreak}
\newcommand{\explique}[2]{\par\smallskip\noindent
  \begin{minipage}[t]{0.6\linewidth}\raggedright #1\end{minipage}\hfill
  \begin{minipage}[t]{0.37\linewidth}\raggedright\leavevmode
    {\small\itshape\color{black!60}#2}\end{minipage}\par}

\setlength{\parskip}{2pt}

\makeatother
\entete{Chapitre 4 --- Calcul littéral --- corrigé}{Seconde}
\setlength{\parindent}{0pt}
\setlength{\jot}{5pt}

\let\casligne\cas
\renewcommand{\cas}[1]{\casligne{#1}\par\nopagebreak}

\begin{document}

\entetefiche{Chapitre 4 --- Calcul littéral}{Seconde}

\partie{Réduire, substituer}

\begin{exo}[Réduire une somme]
\explique{Pour réduire une somme, on rapproche les termes semblables (même partie
littérale : les $a$ ensemble, les $a^{2}$ ensemble, les nombres ensemble), avec le
signe placé devant eux, puis on additionne leurs nombres.}{méthode 1 du cours}
\cas{A}
\begin{calculs}
A &= 7a+5a+2a &&\justif{on recopie : trois termes en $a$}\\
A &= (7+5+2)a &&\justif{on additionne les nombres devant $a$}\\
A &= \res{14a} &&
\end{calculs}
\cas{B}
\begin{calculs}
B &= 6+2a+5a &&\justif{on recopie}\\
B &= 6+7a &&\justif{$2a+5a=7a$ ; le $6$ n'a pas de terme semblable}\\
B &= \res{7a+6} &&\justif{on écrit d'habitude les termes en $a$ d'abord}
\end{calculs}
\rem{Erreur fréquente : écrire $B=13a$. Le nombre $6$ et les $a$ ne sont pas semblables :
$6+7a$ reste une somme de deux termes.}
\cas{C}
\begin{calculs}
C &= 4a^{2}+3+7a^{2} &&\justif{on recopie}\\
C &= 4a^{2}+7a^{2}+3 &&\justif{on rapproche les $a^{2}$}\\
C &= \res{11a^{2}+3} &&
\end{calculs}
\cas{D}
\begin{calculs}
D &= 6a+3a^{2}+4a^{2} &&\justif{on recopie}\\
D &= 3a^{2}+4a^{2}+6a &&\justif{les $a^{2}$ ensemble ; $6a$ n'est pas semblable à $a^{2}$}\\
D &= \res{7a^{2}+6a} &&\justif{on s'arrête : $7a^{2}$ et $6a$ ne se regroupent pas}
\end{calculs}
\cas{E}
\begin{calculs}
E &= 9m-12m+5m &&\justif{on recopie}\\
E &= (9-12+5)m &&\justif{de gauche à droite : $9-12=-3$, puis $-3+5=2$}\\
E &= \res{2m} &&
\end{calculs}
\cas{F}
\begin{calculs}
F &= -6m-3m-4m &&\justif{trois termes négatifs}\\
F &= (-6-3-4)m &&\justif{on ajoute les distances à zéro : $6+3+4=13$}\\
F &= \res{-13m} &&\justif{on garde le signe moins}
\end{calculs}
\cas{G}
\begin{calculs}
G &= 3m-8m^{2}+5m^{2} &&\justif{on recopie}\\
G &= -8m^{2}+5m^{2}+3m &&\justif{les $m^{2}$ avec leur signe}\\
G &= \res{-3m^{2}+3m} &&\justif{$-8+5=-3$}
\end{calculs}
\cas{H}
\begin{calculs}
H &= -6m^{2}+4m^{2}-5m^{2} &&\justif{trois termes semblables}\\
H &= (-6+4-5)m^{2} &&\justif{$-6+4=-2$, puis $-2-5=-7$}\\
H &= \res{-7m^{2}} &&
\end{calculs}
\cas{K}
\begin{calculs}
K &= 3n^{2}+6n-5+4n^{2}+3+8 &&\justif{on recopie}\\
K &= 3n^{2}+4n^{2}+6n-5+3+8 &&\justif{les $n^{2}$, puis les $n$, puis les nombres}\\
K &= \res{7n^{2}+6n+6} &&\justif{$-5+3+8=6$}
\end{calculs}
\cas{L}
\begin{calculs}
L &= 4+5n-3n^{2}+7n^{2}+2n-4 &&\justif{on recopie}\\
L &= -3n^{2}+7n^{2}+5n+2n+4-4 &&\justif{on rapproche les termes semblables}\\
L &= \res{4n^{2}+7n} &&\justif{$4-4=0$ : il n'y a plus de nombre seul}
\end{calculs}
\cas{M}
\begin{calculs}
M &= -7n^{2}+6-4n+3n^{2}-3n-9 &&\justif{on recopie}\\
M &= -7n^{2}+3n^{2}-4n-3n+6-9 &&\justif{chaque terme garde le signe écrit devant lui}\\
M &= \res{-4n^{2}-7n-3} &&\justif{$-7+3=-4$ ; $-4-3=-7$ ; $6-9=-3$}
\end{calculs}
\rem{Contrôle rapide : on remplace la lettre par $1$ dans l'expression de départ et dans le
résultat. Pour $M$ : $-7+6-4+3-3-9=-14$ et $-4-7-3=-14$.}
\end{exo}

\begin{exo}[Réduire un produit]
\explique{Pour réduire un produit, on multiplie les nombres entre eux (règle des
signes), puis les lettres entre elles : $x\times x=x^{2}$.}{un produit se réduit
toujours}
\cas{A}
\begin{calculs}
A &= 6\times 4x &&\justif{on recopie}\\
A &= 6\times 4\times x &&\justif{les nombres, puis la lettre}\\
A &= \res{24x} &&
\end{calculs}
\cas{B}
\begin{calculs}
B &= 3x\times 5x &&\justif{on recopie}\\
B &= 3\times 5\times x\times x &&\justif{les nombres, puis les lettres}\\
B &= \res{15x^{2}} &&\justif{$x\times x=x^{2}$}
\end{calculs}
\cas{C}
\begin{calculs}
C &= 7\times 2x^{2} &&\\
C &= \res{14x^{2}} &&\justif{$7\times 2=14$ ; le $x^{2}$ ne change pas}
\end{calculs}
\cas{D}
\begin{calculs}
D &= 3x^{2}\times 5 &&\\
D &= \res{15x^{2}} &&\justif{$3\times 5=15$}
\end{calculs}
\cas{E}
\begin{calculs}
E &= -4x\times 3x &&\justif{signes contraires : le produit est négatif}\\
E &= -4\times 3\times x\times x &&\\
E &= \res{-12x^{2}} &&
\end{calculs}
\cas{F}
\begin{calculs}
F &= -6\times 5x &&\justif{signes contraires : négatif}\\
F &= \res{-30x} &&\justif{une seule lettre : pas de carré}
\end{calculs}
\cas{G}
\begin{calculs}
G &= 2x(-7x) &&\justif{le $\times$ est sous-entendu devant la parenthèse}\\
G &= 2\times(-7)\times x\times x &&\justif{signes contraires : négatif}\\
G &= \res{-14x^{2}} &&
\end{calculs}
\cas{H}
\begin{calculs}
H &= -3x(-6x) &&\justif{deux nombres négatifs : le produit est positif}\\
H &= (-3)\times(-6)\times x\times x &&\\
H &= \res{18x^{2}} &&
\end{calculs}
\cas{K}
\begin{calculs}
K &= 4(-3x^{2}) &&\justif{signes contraires : négatif}\\
K &= \res{-12x^{2}} &&
\end{calculs}
\cas{L}
\begin{calculs}
L &= -5(-2x^{2}) &&\justif{même signe : positif}\\
L &= \res{10x^{2}} &&
\end{calculs}
\cas{M}
\begin{calculs}
M &= 3x^{2}\times(-4) &&\justif{signes contraires : négatif}\\
M &= \res{-12x^{2}} &&
\end{calculs}
\rem{Ne pas confondre somme et produit : $3x+5x=8x$, mais $3x\times 5x=15x^{2}$. Dans un
produit, les exposants des $x$ s'ajoutent : $x\times x^{2}=x^{3}$.}
\end{exo}

\begin{exo}[Substituer]
\explique{Substituer, c'est remplacer la lettre par sa valeur : on écrit la valeur entre
parenthèses, on rétablit les $\times$, puis on calcule avec les priorités (le carré
d'abord, puis les produits, enfin les sommes).}{méthode 2 du cours}
\cas{1) Pour $x=2$}
\begin{calculs}
B &= 3\times 2^{2}-4\times 2-4 &&\justif{on remplace $x$ par $2$}\\
B &= 3\times 4-8-4 &&\justif{le carré d'abord : $2^{2}=4$}\\
B &= 12-8-4 &&\\
B &= \res{0} &&
\end{calculs}
\cas{1) Pour $x=3$}
\begin{calculs}
B &= 3\times 3^{2}-4\times 3-4 &&\\
B &= 3\times 9-12-4 &&\justif{$3^{2}=9$}\\
B &= 27-12-4 &&\\
B &= \res{11} &&
\end{calculs}
\cas{1) Pour $x=-1$}
\begin{calculs}
B &= 3\times(-1)^{2}-4\times(-1)-4 &&\justif{la valeur négative entre parenthèses}\\
B &= 3\times 1+4-4 &&\justif{$(-1)^{2}=1$ ; $-4\times(-1)=+4$}\\
B &= \res{3} &&
\end{calculs}
\cas{1) Pour $x=-2$}
\begin{calculs}
B &= 3\times(-2)^{2}-4\times(-2)-4 &&\\
B &= 3\times 4+8-4 &&\justif{$(-2)^{2}=4$ : un carré est positif}\\
B &= 12+8-4 &&\\
B &= \res{16} &&
\end{calculs}
\cas{1) Pour $x=\frac13$}
\begin{calculs}
B &= 3\times\left(\tfrac13\right)^{2}-4\times\tfrac13-4 &&\\
B &= 3\times\tfrac19-\tfrac43-4 &&\justif{$\left(\frac13\right)^{2}=\frac{1\times 1}{3\times 3}=\frac19$}\\
B &= \tfrac13-\tfrac43-4 &&\justif{$3\times\frac19=\frac39=\frac13$}\\
B &= -1-4 &&\justif{$\frac13-\frac43=-\frac33=-1$}\\
B &= \res{-5} &&
\end{calculs}
\cas{2) Pour $m=3$ et $p=-2$}
\begin{calculs}
C &= 2\times 3-4\times(-2)+5 &&\justif{chaque lettre remplacée par sa valeur}\\
C &= 6+8+5 &&\justif{$-4\times(-2)=+8$}\\
C &= \res{19} &&
\end{calculs}
\cas{2) Pour $m=-1$ et $p=4$}
\begin{calculs}
C &= 2\times(-1)-4\times 4+5 &&\\
C &= -2-16+5 &&\\
C &= \res{-13} &&\justif{$-18+5=-13$}
\end{calculs}
\rem{Erreur fréquente : écrire $-4\times-2$ sans parenthèses, puis se tromper de signe.
Autre piège : $3\times(-2)^{2}$ vaut $3\times 4=12$ ; le carré se calcule avant le
produit.}
\end{exo}

\begin{exo}[Égales ou non ?]
\cas{1)}
\begin{calculs}
A &= (0-2)(0+5) &&\justif{pour $x=0$, dans chaque expression}\\
A &= -2\times 5 &&\\
A &= -10 &&\\
B &= 2\times 0^{2}-10 &&\\
B &= -10 &&
\end{calculs}
\begin{calculs}
A &= (1-2)(1+5) &&\justif{pour $x=1$}\\
A &= -1\times 6 &&\\
A &= -6 &&\\
B &= 2\times 1^{2}-10 &&\\
B &= -8 &&
\end{calculs}
\explique{Pour $x=0$, $A=\res{-10}$ et $B=\res{-10}$ ; pour $x=1$, $A=\res{-6}$ et
$B=\res{-8}$.}{}
\cas{2)}
\explique{Pour $x=1$, $A$ et $B$ prennent des valeurs différentes : elles ne sont
\res{pas égales}.}{un seul contre-exemple suffit}
\rem{Le test avec $x=0$ a réussi, mais il ne prouve rien : deux expressions égales
doivent donner le même résultat pour \emph{toutes} les valeurs de $x$.}
\cas{3)}
\begin{calculs}
A &= (x-2)(x+5) &&\justif{on recopie}\\
A &= x^{2}+5x-2x-10 &&\justif{chaque terme par chaque terme : $x\times x$ ; $x\times 5$ ; $-2\times x$ ; $-2\times 5$}\\
A &= \res{x^{2}+3x-10} &&
\end{calculs}
\explique{La forme développée de $A$ n'est pas $2x^{2}-10$ : on confirme que $A$ et $B$
ne sont pas égales.}{pour prouver une égalité, on développe}
\end{exo}

\partie{Développer}

\begin{exo}[Distributivité simple]
\explique{Le facteur devant (ou derrière) la parenthèse multiplie \emph{chaque} terme :
$k(a+b)=ka+kb$. On réduit ensuite.}{propriété 2 et méthode 3 du cours}
\cas{A}
\begin{calculs}
A &= 4(3x+5) &&\justif{on recopie}\\
A &= 4\times 3x+4\times 5 &&\justif{$4$ multiplie $3x$, puis $5$}\\
A &= \res{12x+20} &&
\end{calculs}
\cas{B}
\begin{calculs}
B &= (6+3x)\times 2 &&\justif{le facteur est à droite : il multiplie aussi chaque terme}\\
B &= 6\times 2+3x\times 2 &&\\
B &= 12+6x &&\\
B &= \res{6x+12} &&
\end{calculs}
\cas{C}
\begin{calculs}
C &= x(4x+5) &&\justif{le facteur est $x$}\\
C &= x\times 4x+x\times 5 &&\\
C &= \res{4x^{2}+5x} &&\justif{$x\times 4x=4x^{2}$}
\end{calculs}
\cas{D}
\begin{calculs}
D &= 3x(2x+7) &&\\
D &= 3x\times 2x+3x\times 7 &&\\
D &= \res{6x^{2}+21x} &&\justif{$3x\times 2x=6x^{2}$ et $3x\times 7=21x$}
\end{calculs}
\cas{E}
\begin{calculs}
E &= 5x+3(2x+4) &&\justif{le produit d'abord ; le $5x$ attend}\\
E &= 5x+6x+12 &&\\
E &= \res{11x+12} &&
\end{calculs}
\cas{F}
\begin{calculs}
F &= 2(5x+6)+4(2x+3) &&\justif{deux produits, développés chacun}\\
F &= 10x+12+8x+12 &&\\
F &= \res{18x+24} &&
\end{calculs}
\cas{G}
\begin{calculs}
G &= 5(3+2x)+4x &&\\
G &= 15+10x+4x &&\\
G &= \res{14x+15} &&
\end{calculs}
\cas{H}
\begin{calculs}
H &= 3(4x+2)+2(6+3x) &&\\
H &= 12x+6+12+6x &&\\
H &= \res{18x+18} &&
\end{calculs}
\cas{K}
\begin{calculs}
K &= x(3x+5)+2(5x+1) &&\\
K &= 3x^{2}+5x+10x+2 &&\\
K &= \res{3x^{2}+15x+2} &&
\end{calculs}
\cas{L}
\begin{calculs}
L &= 2x(4x+3)+3x(2x+5) &&\\
L &= 8x^{2}+6x+6x^{2}+15x &&\justif{$2x\times 4x=8x^{2}$ ; $3x\times 2x=6x^{2}$}\\
L &= \res{14x^{2}+21x} &&
\end{calculs}
\rem{Erreur fréquente : ne multiplier que le premier terme, $4(3x+5)=12x+5$. Contrôle avec
$x=1$ : $4\times 8=32$ et $12+20=32$.}
\end{exo}

\begin{exo}[Avec des signes moins]
\explique{Le facteur garde son signe et multiplie chaque terme : on applique la règle des
signes à chaque produit. Un $-$ devant une parenthèse change le signe de chaque
terme.}{propriété 2 du cours}
\cas{A}
\begin{calculs}
A &= 6(a+3)-4a &&\\
A &= 6a+18-4a &&\\
A &= \res{2a+18} &&
\end{calculs}
\cas{B}
\begin{calculs}
B &= -7a+3(2a-4) &&\\
B &= -7a+6a-12 &&\justif{$3\times(-4)=-12$}\\
B &= \res{-a-12} &&\justif{$-7a+6a=-1a$, qu'on écrit $-a$}
\end{calculs}
\cas{C}
\begin{calculs}
C &= 4(3+2a)-3(5+2a) &&\justif{le second facteur est $-3$}\\
C &= 12+8a-15-6a &&\justif{$-3\times 5=-15$ et $-3\times 2a=-6a$}\\
C &= \res{2a-3} &&
\end{calculs}
\cas{D}
\begin{calculs}
D &= -5(-3a+2)+2(4a-3) &&\\
D &= 15a-10+8a-6 &&\justif{$-5\times(-3a)=+15a$ ; $-5\times 2=-10$}\\
D &= \res{23a-16} &&
\end{calculs}
\cas{E}
\begin{calculs}
E &= 3b+2b(5+7b) &&\\
E &= 3b+10b+14b^{2} &&\justif{$2b\times 7b=14b^{2}$}\\
E &= \res{14b^{2}+13b} &&
\end{calculs}
\cas{F}
\begin{calculs}
F &= 5b(-3b+2)+4b^{2} &&\\
F &= -15b^{2}+10b+4b^{2} &&\justif{$5b\times(-3b)=-15b^{2}$}\\
F &= \res{-11b^{2}+10b} &&
\end{calculs}
\cas{G}
\begin{calculs}
G &= 4b(-3b-2)-b(2b-9) &&\justif{le second facteur est $-b$}\\
G &= -12b^{2}-8b-2b^{2}+9b &&\justif{$-b\times 2b=-2b^{2}$ et $-b\times(-9)=+9b$}\\
G &= \res{-14b^{2}+b} &&
\end{calculs}
\cas{H}
\begin{calculs}
H &= -3b(4b-2)+6b(-2b-3) &&\\
H &= -12b^{2}+6b-12b^{2}-18b &&\justif{$-3b\times(-2)=+6b$}\\
H &= \res{-24b^{2}-12b} &&
\end{calculs}
\rem{Erreur fréquente dans $C$ et $G$ : ne changer que le signe du premier terme. Le
facteur $-3$ (ou $-b$) multiplie \emph{chaque} terme de la parenthèse.}
\end{exo}

\begin{exo}[Double distributivité]
\explique{Chaque terme de la première parenthèse multiplie chaque terme de la seconde :
$(a+b)(c+d)=ac+ad+bc+bd$, soit quatre produits, chacun avec son signe.}{méthode 3 du
cours, exemple e)}
\cas{A}
\begin{calculs}
A &= (2x+5)(4x+3) &&\justif{on recopie}\\
A &= 8x^{2}+6x+20x+15 &&\justif{$2x\times 4x$ ; $2x\times 3$ ; $5\times 4x$ ; $5\times 3$}\\
A &= \res{8x^{2}+26x+15} &&
\end{calculs}
\cas{B}
\begin{calculs}
B &= (3+4x)(2x+7) &&\\
B &= 6x+21+8x^{2}+28x &&\\
B &= \res{8x^{2}+34x+21} &&\justif{on range : $x^{2}$, puis $x$, puis le nombre}
\end{calculs}
\cas{C}
\begin{calculs}
C &= (3x-4)(-2x+5) &&\\
C &= -6x^{2}+15x+8x-20 &&\justif{$-4\times(-2x)=+8x$ ; $-4\times 5=-20$}\\
C &= \res{-6x^{2}+23x-20} &&
\end{calculs}
\cas{D}
\begin{calculs}
D &= (-5+2x)(3-4x) &&\\
D &= -15+20x+6x-8x^{2} &&\justif{$-5\times(-4x)=+20x$ ; $2x\times(-4x)=-8x^{2}$}\\
D &= \res{-8x^{2}+26x-15} &&
\end{calculs}
\cas{E}
\begin{calculs}
E &= (-4x-2)(-3x-5) &&\\
E &= 12x^{2}+20x+6x+10 &&\justif{quatre produits de deux négatifs : tous positifs}\\
E &= \res{12x^{2}+26x+10} &&
\end{calculs}
\cas{F}
\begin{calculs}
F &= (2-3x)(4-5x) &&\\
F &= 8-10x-12x+15x^{2} &&\justif{$-3x\times(-5x)=+15x^{2}$}\\
F &= \res{15x^{2}-22x+8} &&
\end{calculs}
\cas{G}
\begin{calculs}
G &= (-6x+5)(-3x+2) &&\\
G &= 18x^{2}-12x-15x+10 &&\\
G &= \res{18x^{2}-27x+10} &&
\end{calculs}
\cas{H}
\begin{calculs}
H &= (-3+2x)(4x-5) &&\\
H &= -12x+15+8x^{2}-10x &&\\
H &= \res{8x^{2}-22x+15} &&
\end{calculs}
\rem{Contrôle avec $x=1$ : pour $C$, $(3-4)(-2+5)=-3$ et $-6+23-20=-3$. Pour $F$,
$(2-3)(4-5)=1$ et $15-22+8=1$.}
\end{exo}

\begin{exo}[Une différence de produits]
\explique{Le $-$ devant un produit porte sur tout ce produit : on développe d'abord ce
produit entre crochets, puis le $-$ change le signe de chaque terme du crochet.}{propriété
2 : pour soustraire un produit, on le développe entre crochets}
\cas{A}
\begin{calculs}
A &= (4x+3)(2x+5)-(3x+2)(2x-5) &&\justif{on recopie}\\
A &= 8x^{2}+20x+6x+15-\big[6x^{2}-15x+4x-10\big] &&\justif{chaque produit à part}\\
A &= 8x^{2}+26x+15-\big[6x^{2}-11x-10\big] &&\justif{on réduit dans le crochet}\\
A &= 8x^{2}+26x+15-6x^{2}+11x+10 &&\justif{le $-$ change tous les signes du crochet}\\
A &= \res{2x^{2}+37x+25} &&
\end{calculs}
\cas{B}
\begin{calculs}
B &= (-2x+7)(3x-4)+(5x-2)(2x-3) &&\justif{ici une somme : pas de changement de signe}\\
B &= -6x^{2}+8x+21x-28+10x^{2}-15x-4x+6 &&\\
B &= \res{4x^{2}+10x-22} &&\justif{$-6+10=4$ ; $8+21-15-4=10$ ; $-28+6=-22$}
\end{calculs}
\cas{C}
\begin{calculs}
C &= (-4x-5)(-3x+6)-(2+3x)(5-4x) &&\\
C &= 12x^{2}-24x+15x-30-\big[10-8x+15x-12x^{2}\big] &&\\
C &= 12x^{2}-9x-30-\big[-12x^{2}+7x+10\big] &&\\
C &= 12x^{2}-9x-30+12x^{2}-7x-10 &&\justif{le $-$ change tous les signes}\\
C &= \res{24x^{2}-16x-40} &&
\end{calculs}
\cas{D}
\begin{calculs}
D &= (6x+1)(3x-2)-(5x-2)(4-3x) &&\\
D &= 18x^{2}-12x+3x-2-\big[20x-15x^{2}-8+6x\big] &&\\
D &= 18x^{2}-9x-2-\big[-15x^{2}+26x-8\big] &&\\
D &= 18x^{2}-9x-2+15x^{2}-26x+8 &&\\
D &= \res{33x^{2}-35x+6} &&
\end{calculs}
\rem{Erreur fréquente : écrire $-(3x+2)(2x-5)=-6x^{2}-15x+\ldots$, en ne changeant que le
premier signe. Le crochet évite cette erreur.}
\end{exo}

\begin{exo}[Crochets et trois facteurs]
\cas{A}
\begin{calculs}
A &= (6-3x)-4-(5x-2) &&\justif{un $+$ devant la première parenthèse : on l'enlève}\\
A &= 6-3x-4-5x+2 &&\justif{le $-$ change les signes de $5x-2$}\\
A &= \res{4-8x} &&\justif{$6-4+2=4$ et $-3x-5x=-8x$}
\end{calculs}
\cas{B}
\begin{calculs}
B &= (4x-3)+5x-(6x+2) &&\\
B &= 4x-3+5x-6x-2 &&\\
B &= \res{3x-5} &&
\end{calculs}
\cas{C}
\explique{Des parenthèses dans des crochets : on commence par les parenthèses les plus
intérieures.}{comme pour les priorités de calcul}
\begin{calculs}
C &= 4-\big[6-(2x+3)\big]+(3-4x) &&\justif{on recopie}\\
C &= 4-\big[6-2x-3\big]+3-4x &&\justif{la parenthèse intérieure}\\
C &= 4-\big[3-2x\big]+3-4x &&\justif{on réduit le crochet}\\
C &= 4-3+2x+3-4x &&\justif{le $-$ change les signes du crochet}\\
C &= \res{4-2x} &&
\end{calculs}
\cas{D}
\begin{calculs}
D &= 5x+\big[2-(6-x)\big]-\big[3x-(x-4)\big] &&\\
D &= 5x+\big[2-6+x\big]-\big[3x-x+4\big] &&\justif{les parenthèses intérieures}\\
D &= 5x+\big[x-4\big]-\big[2x+4\big] &&\\
D &= 5x+x-4-2x-4 &&\\
D &= \res{4x-8} &&
\end{calculs}
\cas{E}
\begin{calculs}
E &= 3x(2x+5)-4(2x+5)(x-3) &&\justif{trois facteurs : d'abord les deux parenthèses}\\
E &= 6x^{2}+15x-4\big[2x^{2}-6x+5x-15\big] &&\\
E &= 6x^{2}+15x-4\big[2x^{2}-x-15\big] &&\\
E &= 6x^{2}+15x-8x^{2}+4x+60 &&\justif{puis $-4$ multiplie chaque terme du crochet}\\
E &= \res{-2x^{2}+19x+60} &&
\end{calculs}
\cas{F}
\begin{calculs}
F &= (x+2)(3x-4)(-x+3) &&\justif{deux facteurs d'abord}\\
F &= \big[3x^{2}-4x+6x-8\big](-x+3) &&\\
F &= (3x^{2}+2x-8)(-x+3) &&\justif{puis chaque terme par chaque terme : six produits}\\
F &= -3x^{3}+9x^{2}-2x^{2}+6x+8x-24 &&\justif{$3x^{2}\times(-x)=-3x^{3}$}\\
F &= \res{-3x^{3}+7x^{2}+14x-24} &&
\end{calculs}
\rem{Dans $E$, on ne multiplie pas $-4$ par $(2x+5)$ et par $(x-3)$ à la fois : le
$-4$ ne multiplie qu'une fois le produit développé. Contrôle de $F$ avec $x=1$ :
$3\times(-1)\times 2=-6$ et $-3+7+14-24=-6$.}
\end{exo}

\partie{Identités remarquables}

\begin{exo}[Le carré d'une somme]
\explique{On reconnaît $(a+b)^{2}=a^{2}+2ab+b^{2}$ : on écrit $a$ et $b$, puis les trois
termes, avec des parenthèses autour de $a$ ou de $b$ quand ils contiennent un
nombre et une lettre.}{méthode 4 du cours}
\cas{A}
\begin{calculs}
A &= (3x+4)^{2} &&\justif{$(a+b)^{2}$ avec $a=3x$ et $b=4$}\\
A &= (3x)^{2}+2\times 3x\times 4+4^{2} &&\justif{carré, double produit, carré}\\
A &= \res{9x^{2}+24x+16} &&\justif{$(3x)^{2}=9x^{2}$}
\end{calculs}
\cas{B}
\begin{calculs}
B &= (5+2x)^{2} &&\justif{$a=5$ et $b=2x$}\\
B &= 5^{2}+2\times 5\times 2x+(2x)^{2} &&\\
B &= \res{25+20x+4x^{2}} &&
\end{calculs}
\cas{C}
\begin{calculs}
C &= (4a+3)^{2} &&\justif{le $a$ de la formule est ici $4a$, et le $b$ est $3$}\\
C &= (4a)^{2}+2\times 4a\times 3+3^{2} &&\\
C &= \res{16a^{2}+24a+9} &&
\end{calculs}
\cas{D}
\begin{calculs}
D &= (3+5b)^{2} &&\justif{le $a$ de la formule est ici $3$, et le $b$ est $5b$}\\
D &= 3^{2}+2\times 3\times 5b+(5b)^{2} &&\\
D &= \res{9+30b+25b^{2}} &&
\end{calculs}
\cas{E}
\begin{calculs}
E &= (x+6)^{2} &&\justif{$a=x$ et $b=6$}\\
E &= x^{2}+2\times x\times 6+6^{2} &&\\
E &= \res{x^{2}+12x+36} &&
\end{calculs}
\cas{F}
\begin{calculs}
F &= (4x+1)^{2} &&\justif{$a=4x$ et $b=1$}\\
F &= (4x)^{2}+2\times 4x\times 1+1^{2} &&\\
F &= \res{16x^{2}+8x+1} &&
\end{calculs}
\cas{G}
\begin{calculs}
G &= (3m+5)^{2} &&\justif{$a=3m$ et $b=5$}\\
G &= (3m)^{2}+2\times 3m\times 5+5^{2} &&\\
G &= \res{9m^{2}+30m+25} &&
\end{calculs}
\cas{H}
\begin{calculs}
H &= (4+3r)^{2} &&\justif{$a=4$ et $b=3r$}\\
H &= 4^{2}+2\times 4\times 3r+(3r)^{2} &&\\
H &= \res{16+24r+9r^{2}} &&
\end{calculs}
\rem{Erreur la plus fréquente : $(3x+4)^{2}=9x^{2}+16$, sans le double produit. Contrôle
avec $x=1$ : $(3+4)^{2}=49$ et $9+24+16=49$, alors que $9+16=25$. Autre erreur :
$(3x)^{2}=3x^{2}$ ; le carré porte sur le $3$ aussi.}
\end{exo}

\begin{exo}[Le carré d'une différence]
\explique{On reconnaît $(a-b)^{2}=a^{2}-2ab+b^{2}$ : seul le double produit est
précédé d'un $-$ ; les deux carrés sont positifs.}{méthode 4 du cours}
\cas{A}
\begin{calculs}
A &= (4x-3)^{2} &&\justif{$(a-b)^{2}$ avec $a=4x$ et $b=3$}\\
A &= (4x)^{2}-2\times 4x\times 3+3^{2} &&\\
A &= \res{16x^{2}-24x+9} &&
\end{calculs}
\cas{B}
\begin{calculs}
B &= (3x-5)^{2} &&\justif{$a=3x$ et $b=5$}\\
B &= (3x)^{2}-2\times 3x\times 5+5^{2} &&\\
B &= \res{9x^{2}-30x+25} &&
\end{calculs}
\cas{C}
\begin{calculs}
C &= (4-3x)^{2} &&\justif{$a=4$ et $b=3x$}\\
C &= 4^{2}-2\times 4\times 3x+(3x)^{2} &&\\
C &= \res{16-24x+9x^{2}} &&
\end{calculs}
\cas{D}
\begin{calculs}
D &= (5-2x)^{2} &&\justif{$a=5$ et $b=2x$}\\
D &= 5^{2}-2\times 5\times 2x+(2x)^{2} &&\\
D &= \res{25-20x+4x^{2}} &&
\end{calculs}
\cas{E}
\begin{calculs}
E &= (5a-3)^{2} &&\justif{le $a$ de la formule est ici $5a$, et le $b$ est $3$}\\
E &= (5a)^{2}-2\times 5a\times 3+3^{2} &&\\
E &= \res{25a^{2}-30a+9} &&
\end{calculs}
\cas{F}
\begin{calculs}
F &= (2-5b)^{2} &&\justif{le $a$ de la formule est ici $2$, et le $b$ est $5b$}\\
F &= 2^{2}-2\times 2\times 5b+(5b)^{2} &&\\
F &= \res{4-20b+25b^{2}} &&
\end{calculs}
\cas{G}
\begin{calculs}
G &= (3m-1)^{2} &&\justif{$a=3m$ et $b=1$}\\
G &= (3m)^{2}-2\times 3m\times 1+1^{2} &&\\
G &= \res{9m^{2}-6m+1} &&
\end{calculs}
\cas{H}
\begin{calculs}
H &= (n-4)^{2} &&\justif{$a=n$ et $b=4$}\\
H &= n^{2}-2\times n\times 4+4^{2} &&\\
H &= \res{n^{2}-8n+16} &&
\end{calculs}
\rem{Erreur fréquente : $(4x-3)^{2}=16x^{2}-24x-9$. Le dernier terme est $b^{2}=3^{2}=9$,
positif. Contrôle avec $x=1$ : $(4-3)^{2}=1$ et $16-24+9=1$.}
\end{exo}

\begin{exo}[Somme par différence]
\explique{On reconnaît $(a+b)(a-b)=a^{2}-b^{2}$ : les deux parenthèses ont les
\emph{mêmes} termes, l'une avec $+$, l'autre avec $-$.}{avant d'appliquer la formule, on
vérifie que les termes sont les mêmes}
\cas{A}
\begin{calculs}
A &= (4x+3)(4x-3) &&\justif{$a=4x$ et $b=3$}\\
A &= (4x)^{2}-3^{2} &&\\
A &= \res{16x^{2}-9} &&
\end{calculs}
\cas{B}
\explique{$2x+5$ et $2x-3$ n'ont pas les mêmes nombres ($5$ et $3$) : ce n'est pas une
identité remarquable. On utilise la double distributivité.}{le produit piège annoncé par
l'énoncé}
\begin{calculs}
B &= (2x+5)(2x-3) &&\\
B &= 4x^{2}-6x+10x-15 &&\justif{quatre produits}\\
B &= \res{4x^{2}+4x-15} &&
\end{calculs}
\cas{C}
\begin{calculs}
C &= (3+5x)(3-5x) &&\justif{$a=3$ et $b=5x$}\\
C &= 3^{2}-(5x)^{2} &&\\
C &= \res{9-25x^{2}} &&
\end{calculs}
\cas{D}
\begin{calculs}
D &= (3a+4)(3a-4) &&\justif{$a=3a$ et $b=4$}\\
D &= (3a)^{2}-4^{2} &&\\
D &= \res{9a^{2}-16} &&
\end{calculs}
\cas{E}
\begin{calculs}
E &= (b-5)(b+5) &&\justif{l'ordre des parenthèses ne change rien : $a=b$ et $b=5$}\\
E &= b^{2}-5^{2} &&\\
E &= \res{b^{2}-25} &&
\end{calculs}
\cas{F}
\begin{calculs}
F &= (7+c)(7-c) &&\justif{$a=7$ et $b=c$}\\
F &= 7^{2}-c^{2} &&\\
F &= \res{49-c^{2}} &&
\end{calculs}
\rem{Le résultat est toujours le carré du terme qui garde son signe, moins le carré du
terme qui change de signe : dans $C$, c'est $9-25x^{2}$, et non $25x^{2}-9$.}
\end{exo}

\begin{exo}[Choisir la bonne formule]
\explique{Avant de développer, on regarde la forme : un nombre devant une parenthèse
(distributivité), un carré (première ou deuxième identité), deux parenthèses aux mêmes
termes (troisième identité), sinon la double distributivité.}{on choisit la formule la
plus rapide}
\cas{A}
\begin{calculs}
A &= 3(2x+5) &&\justif{distributivité simple}\\
A &= \res{6x+15} &&
\end{calculs}
\cas{B}
\begin{calculs}
B &= (2x+5)^{2} &&\justif{$(a+b)^{2}$ avec $a=2x$ et $b=5$}\\
B &= (2x)^{2}+2\times 2x\times 5+5^{2} &&\\
B &= \res{4x^{2}+20x+25} &&
\end{calculs}
\cas{C}
\begin{calculs}
C &= (3x-4)^{2} &&\justif{$(a-b)^{2}$ avec $a=3x$ et $b=4$}\\
C &= (3x)^{2}-2\times 3x\times 4+4^{2} &&\\
C &= \res{9x^{2}-24x+16} &&
\end{calculs}
\cas{D}
\begin{calculs}
D &= (5x-3)(5x+3) &&\justif{$(a-b)(a+b)$ avec $a=5x$ et $b=3$}\\
D &= (5x)^{2}-3^{2} &&\\
D &= \res{25x^{2}-9} &&
\end{calculs}
\cas{E}
\begin{calculs}
E &= (7x+2)(3x-5) &&\justif{pas les mêmes termes : double distributivité}\\
E &= 21x^{2}-35x+6x-10 &&\\
E &= \res{21x^{2}-29x-10} &&
\end{calculs}
\cas{F}
\begin{calculs}
F &= (3+7x)(3-7x) &&\justif{$a=3$ et $b=7x$}\\
F &= 3^{2}-(7x)^{2} &&\\
F &= \res{9-49x^{2}} &&
\end{calculs}
\cas{G}
\begin{calculs}
G &= (2x+3)^{2}-(x-4)^{2} &&\justif{deux identités ; le second carré se soustrait}\\
G &= 4x^{2}+12x+9-\big[x^{2}-8x+16\big] &&\justif{le second carré entre crochets}\\
G &= 4x^{2}+12x+9-x^{2}+8x-16 &&\justif{le $-$ change tous les signes}\\
G &= \res{3x^{2}+20x-7} &&
\end{calculs}
\cas{H}
\begin{calculs}
H &= (x-5)(x+5)-(2x+1)(2x-1) &&\justif{deux fois la troisième identité}\\
H &= x^{2}-25-\big[4x^{2}-1\big] &&\\
H &= x^{2}-25-4x^{2}+1 &&\\
H &= \res{-3x^{2}-24} &&
\end{calculs}
\cas{K}
\begin{calculs}
K &= (x+4)^{2}+3(2x-1)(x+2) &&\justif{une identité, puis trois facteurs}\\
K &= x^{2}+8x+16+3\big[2x^{2}+4x-x-2\big] &&\justif{les deux parenthèses d'abord}\\
K &= x^{2}+8x+16+3\big[2x^{2}+3x-2\big] &&\\
K &= x^{2}+8x+16+6x^{2}+9x-6 &&\justif{puis le $3$ sur chaque terme}\\
K &= \res{7x^{2}+17x+10} &&
\end{calculs}
\cas{L}
\begin{calculs}
L &= (3x+1)(3x-1)-2(x-3)^{2} &&\justif{le carré se calcule avant le produit par $2$}\\
L &= 9x^{2}-1-2\big[x^{2}-6x+9\big] &&\\
L &= 9x^{2}-1-2x^{2}+12x-18 &&\justif{$-2$ multiplie chaque terme du crochet}\\
L &= \res{7x^{2}+12x-19} &&
\end{calculs}
\rem{Dans $L$, on ne calcule pas $\big(2(x-3)\big)^{2}$ : le carré ne porte que sur $(x-3)$.
Contrôle de $L$ avec $x=1$ : $4\times 2-2\times 4=0$ et $7+12-19=0$.}
\end{exo}

\partie{Factoriser}

\begin{exo}[Compléter les factorisations]
\explique{Factoriser, c'est faire le chemin inverse du développement : $ka+kb=k(a+b)$. On
cherche ce qui reste de chaque terme quand on enlève le facteur commun.}{définition 2 du
cours}
\cas{a)}
\begin{calculs}
5\times x+5\times 8 &= 5(x+8) &&\justif{le facteur commun $5$ est déjà écrit}
\end{calculs}
\explique{On complète : $\res{5(x+8)}$.}{contrôle : $5(x+8)=5x+40$}
\cas{b)}
\begin{calculs}
12x-8 &= 4\times 3x-4\times 2 &&\justif{$12=4\times 3$ et $8=4\times 2$}\\
12x-8 &= \res{4(3x-2)} &&
\end{calculs}
\cas{c)}
\begin{calculs}
10x-15 &= 5\times 2x-5\times 3 &&\justif{$2x$ devient $10x$ en multipliant par $5$}\\
10x-15 &= \res{5(2x-3)} &&\justif{contrôle : $5\times 3=15$}
\end{calculs}
\cas{d)}
\begin{calculs}
8x^{2}+12x &= 4x\times 2x+4x\times 3 &&\justif{$8x^{2}=4x\times 2x$ et $12x=4x\times 3$}\\
8x^{2}+12x &= \res{4x(2x+3)} &&
\end{calculs}
\cas{e)}
\begin{calculs}
x^{2}-9x &= x\times x-x\times 9 &&\justif{$x^{2}=x\times x$ : $x$ est dans les deux termes}\\
x^{2}-9x &= \res{x(x-9)} &&
\end{calculs}
\cas{f)}
\begin{calculs}
(x+3)(3x-2)+(x+3)(x+4) &= (x+3)\big[(3x-2)+(x+4)\big] &&\justif{le facteur commun est la parenthèse $(x+3)$}\\
(x+3)(3x-2)+(x+3)(x+4) &= \res{(x+3)(4x+2)} &&\justif{on réduit le crochet : $3x+x=4x$ et $-2+4=2$}
\end{calculs}
\end{exo}

\begin{exo}[Un nombre ou $x$ en facteur]
\explique{On cherche le plus grand nombre qui divise tous les nombres (tables de
multiplication à l'envers), et on ajoute $x$ si chaque terme contient $x$.}{on factorise
le plus possible}
\cas{A}
\begin{calculs}
A &= 48a-36b &&\justif{$12$ divise $48$ et $36$ : $48=12\times 4$ et $36=12\times 3$}\\
A &= 12\times 4a-12\times 3b &&\\
A &= \res{12(4a-3b)} &&\justif{$4$ et $3$ n'ont plus de diviseur commun}
\end{calculs}
\cas{B}
\begin{calculs}
B &= 150a-90b &&\justif{$150=30\times 5$ et $90=30\times 3$}\\
B &= 30\times 5a-30\times 3b &&\\
B &= \res{30(5a-3b)} &&
\end{calculs}
\cas{C}
\begin{calculs}
C &= 24x-36 &&\justif{$24=12\times 2$ et $36=12\times 3$}\\
C &= 12\times 2x-12\times 3 &&\\
C &= \res{12(2x-3)} &&
\end{calculs}
\rem{$C=6(4x-6)$ ou $C=4(6x-9)$ sont justes mais incomplets : on peut encore sortir un
facteur de la parenthèse. On prend le plus grand facteur commun, ici $12$.}
\cas{D}
\begin{calculs}
D &= 35x+13x^{2} &&\justif{$35$ et $13$ n'ont pas de diviseur commun ; $x$ est dans les deux termes}\\
D &= x\times 35+x\times 13x &&\\
D &= \res{x(35+13x)} &&
\end{calculs}
\cas{E}
\begin{calculs}
E &= 28x^{2}-7x &&\justif{$28=7\times 4$ ; $x^{2}=x\times x$ : facteur commun $7x$}\\
E &= 7x\times 4x-7x\times 1 &&\justif{$7x=7x\times 1$}\\
E &= \res{7x(4x-1)} &&
\end{calculs}
\rem{Erreur fréquente : $E=7x(4x)$, en oubliant le $1$. Il reste $1$, pas $0$, quand on
enlève $7x$ de $7x$ ; contrôle : $7x(4x-1)=28x^{2}-7x$.}
\cas{F}
\begin{calculs}
F &= 5x^{3}+10x^{2}+15x &&\justif{$5$ divise chaque coefficient ; chaque terme contient $x$}\\
F &= 5x\times x^{2}+5x\times 2x+5x\times 3 &&\justif{$x^{3}=x\times x^{2}$}\\
F &= \res{5x(x^{2}+2x+3)} &&
\end{calculs}
\end{exo}

\begin{exo}[Une parenthèse en facteur]
\explique{Le facteur commun est une parenthèse entière, présente dans les deux termes. On
l'écrit devant, puis on recopie ce qui reste de chaque terme entre crochets, chacun entre
parenthèses et avec son signe.}{méthode 5 du cours, exemples c) et d)}
\cas{A}
\begin{calculs}
A &= (4x+3)(2x-5)+(4x+3)(3x+7) &&\justif{$(4x+3)$ est dans les deux termes}\\
A &= (4x+3)\big[(2x-5)+(3x+7)\big] &&\\
A &= \res{(4x+3)(5x+2)} &&\justif{$2x+3x=5x$ et $-5+7=2$}
\end{calculs}
\cas{B}
\begin{calculs}
B &= (6x-5)(2x+3)+(6x-5)(4x-1) &&\justif{facteur commun $(6x-5)$}\\
B &= (6x-5)\big[(2x+3)+(4x-1)\big] &&\\
B &= \res{(6x-5)(6x+2)} &&
\end{calculs}
\rem{On peut aller plus loin : $6x+2=2(3x+1)$, donc $B=2(6x-5)(3x+1)$.}
\cas{C}
\begin{calculs}
C &= (3x+5)(2x-7)+(2x-7)(4x+1) &&\justif{le facteur commun $(2x-7)$ n'est pas à la même place}\\
C &= (2x-7)\big[(3x+5)+(4x+1)\big] &&\\
C &= \res{(2x-7)(7x+6)} &&
\end{calculs}
\cas{D}
\begin{calculs}
D &= (8r-2)(5r+3)+(3r-4)(5r+3) &&\justif{facteur commun $(5r+3)$}\\
D &= (5r+3)\big[(8r-2)+(3r-4)\big] &&\\
D &= \res{(5r+3)(11r-6)} &&
\end{calculs}
\cas{E}
\begin{calculs}
E &= (x-4)(x+2)-(x-4)(2x+5) &&\justif{facteur commun $(x-4)$ ; le $-$ porte sur tout le second reste}\\
E &= (x-4)\big[(x+2)-(2x+5)\big] &&\\
E &= (x-4)(x+2-2x-5) &&\justif{$-(2x+5)=-2x-5$}\\
E &= \res{(x-4)(-x-3)} &&
\end{calculs}
\rem{Erreur fréquente : écrire $(x-4)(x+2-2x+5)$. Le $-$ change le signe de $2x$ \emph{et}
de $5$. Contrôle avec $x=0$ : $-4\times 2-(-4)\times 5=12$ et $-4\times(-3)=12$.}
\cas{F}
\begin{calculs}
F &= x(2y-5)+3(2y-5) &&\justif{deux lettres, mais une parenthèse commune $(2y-5)$}\\
F &= (2y-5)\big[x+3\big] &&\\
F &= \res{(2y-5)(x+3)} &&
\end{calculs}
\rem{On contrôle chaque factorisation en la redéveloppant : pour $A$,
$(4x+3)(5x+2)=20x^{2}+23x+6$, comme $8x^{2}-14x-15+12x^{2}+37x+21$.}
\end{exo}

\partie{Factoriser avec une identité remarquable}

\begin{exo}[Différence de deux carrés]
\explique{Deux carrés séparés par un $-$ : on écrit chaque terme comme un carré, $a^{2}$
et $b^{2}$, puis $a^{2}-b^{2}=(a+b)(a-b)$.}{méthode 6 du cours}
\cas{A}
\begin{calculs}
A &= 25x^{2}-16 &&\justif{$25x^{2}=(5x)^{2}$ et $16=4^{2}$}\\
A &= (5x)^{2}-4^{2} &&\justif{$a=5x$ et $b=4$}\\
A &= \res{(5x+4)(5x-4)} &&
\end{calculs}
\cas{B}
\begin{calculs}
B &= 1-9x^{2} &&\justif{$1=1^{2}$ et $9x^{2}=(3x)^{2}$}\\
B &= 1^{2}-(3x)^{2} &&\justif{$a=1$ et $b=3x$}\\
B &= \res{(1+3x)(1-3x)} &&
\end{calculs}
\cas{C}
\explique{$C=4x^{2}+9=(2x)^{2}+3^{2}$ est une \emph{somme} de deux carrés : elle
\res{ne se factorise pas} avec les identités remarquables.}{il faut un signe $-$ entre
les carrés}
\cas{D}
\begin{calculs}
D &= 36R^{2}-25 &&\justif{$a=6R$ et $b=5$}\\
D &= (6R)^{2}-5^{2} &&\\
D &= \res{(6R+5)(6R-5)} &&
\end{calculs}
\cas{E}
\explique{$5x^{2}$ et $5$ ne sont pas des carrés « ronds », mais ils ont le facteur
commun $5$ : on le sort d'abord.}{on regarde toujours d'abord s'il y a un facteur commun}
\begin{calculs}
E &= 5x^{2}-5 &&\\
E &= 5(x^{2}-1) &&\justif{facteur commun $5$}\\
E &= 5(x^{2}-1^{2}) &&\justif{dans la parenthèse : $a=x$ et $b=1$}\\
E &= \res{5(x+1)(x-1)} &&
\end{calculs}
\cas{F}
\begin{calculs}
F &= -x^{2}+9 &&\justif{on change l'ordre pour voir la différence}\\
F &= 9-x^{2} &&\\
F &= 3^{2}-x^{2} &&\justif{$a=3$ et $b=x$}\\
F &= \res{(3+x)(3-x)} &&
\end{calculs}
\rem{Contrôle en redéveloppant : $(3+x)(3-x)=9-x^{2}$. Dans $F$, ne pas écrire
$(x+3)(x-3)$ : cela vaut $x^{2}-9$, l'opposé.}
\end{exo}

\begin{exo}[Carrés parfaits]
\explique{Trois termes dont deux carrés : on trouve $a$ et $b$ (les racines des deux
carrés), puis on contrôle que le troisième terme est bien $2ab$ ; son signe donne
$(a+b)^{2}$ ou $(a-b)^{2}$.}{méthode 6 du cours}
\cas{A}
\explique{$9x^{2}=(3x)^{2}$ et $4=2^{2}$, donc $a=3x$ et $b=2$.}{on trouve $a$ et $b$}
\explique{$2ab=2\times 3x\times 2=12x$ : c'est bien le terme restant, avec un $+$. Donc
$A=\res{(3x+2)^{2}}$.}{on contrôle le double produit}
\cas{B}
\explique{$25a^{2}=(5a)^{2}$ et $9=3^{2}$ ; le double produit serait
$2\times 5a\times 3=30a$. Or le terme restant est $15a$ : $B$ \res{ne se factorise pas}
avec une identité.}{le contrôle échoue}
\cas{C}
\explique{On range d'abord : $C=16d^{2}+24d+9$. Puis $16d^{2}=(4d)^{2}$ et $9=3^{2}$ ;
$2\times 4d\times 3=24d$ est le terme restant. Donc $C=\res{(4d+3)^{2}}$.}{les carrés ne
sont pas toujours aux extrémités}
\cas{D}
\explique{$36=6^{2}$ et $x^{2}=x^{2}$ ; $2\times 6\times x=12x$ et le terme restant est
$-12x$. Donc $D=\res{(6-x)^{2}}$.}{le signe $-$ donne $(a-b)^{2}$}
\rem{$(x-6)^{2}$ est juste aussi : $(6-x)^{2}$ et $(x-6)^{2}$ sont égaux, car deux nombres
opposés ont le même carré.}
\cas{E}
\explique{$E=9x^{2}+6x-1$ : le dernier terme $-1$ est négatif, ce n'est pas un carré.
$E$ \res{ne se factorise pas} avec une identité.}{dans $a^{2}\pm 2ab+b^{2}$, les deux
carrés sont précédés d'un $+$}
\cas{F}
\explique{On range : $F=16b^{2}-40b+25$. Puis $16b^{2}=(4b)^{2}$ et $25=5^{2}$ ;
$2\times 4b\times 5=40b$ et le terme restant est $-40b$. Donc $F=\res{(4b-5)^{2}}$.}{}
\rem{Contrôle de $A$ en redéveloppant : $(3x+2)^{2}=9x^{2}+12x+4$.}
\end{exo}

\begin{exo}[Quelle méthode ?]
\explique{Pour chaque expression, on se pose les questions dans l'ordre : un facteur
commun ? deux carrés séparés par un $-$ ? trois termes dont deux carrés ?}{méthodes 5 et 6
du cours}
\cas{A}
\begin{calculs}
A &= 49x^{2}-16 &&\justif{deux carrés séparés par un $-$ : $a=7x$ et $b=4$}\\
A &= (7x)^{2}-4^{2} &&\\
A &= \res{(7x+4)(7x-4)} &&
\end{calculs}
\cas{B}
\begin{calculs}
B &= 16x^{2}+7x &&\justif{$7x$ n'est pas un carré, mais $x$ est commun}\\
B &= x\times 16x+x\times 7 &&\\
B &= \res{x(16x+7)} &&
\end{calculs}
\cas{C}
\explique{$4x^{2}=(2x)^{2}$ et $25=5^{2}$ ; $2\times 2x\times 5=20x$ est le terme restant.
Donc $C=\res{(2x+5)^{2}}$.}{trois termes : carré parfait}
\cas{D}
\explique{$64=8^{2}$ et $9x^{2}=(3x)^{2}$ ; $2\times 8\times 3x=48x$ et le terme restant
est $-48x$. Donc $D=\res{(8-3x)^{2}}$.}{signe $-$ : $(a-b)^{2}$}
\cas{E}
\begin{calculs}
E &= 5a-4a^{2} &&\justif{$5a$ n'est pas un carré ; $a$ est commun}\\
E &= a\times 5-a\times 4a &&\\
E &= \res{a(5-4a)} &&
\end{calculs}
\cas{F}
\explique{On range : $F=25f^{2}-10f+1$. Puis $25f^{2}=(5f)^{2}$ et $1=1^{2}$ ;
$2\times 5f\times 1=10f$ et le terme restant est $-10f$. Donc
$F=\res{(5f-1)^{2}}$.}{ou $(1-5f)^{2}$, qui est égal}
\end{exo}

\partie{Factorisations avancées}

\begin{exo}[Identité remarquable sur des expressions]
\explique{$a$ et $b$ sont ici des expressions : on les écrit entre parenthèses dans les
crochets de $(a+b)(a-b)$, puis on réduit chaque crochet.}{propriété 5 et méthode 7 du
cours, exemple a)}
\cas{A}
\begin{calculs}
A &= (7x+5)^{2}-(3x+1)^{2} &&\justif{$a=7x+5$ et $b=3x+1$}\\
A &= \big[(7x+5)+(3x+1)\big]\big[(7x+5)-(3x+1)\big] &&\\
A &= (7x+5+3x+1)(7x+5-3x-1) &&\justif{le $-$ change les signes de $3x+1$}\\
A &= \res{(10x+6)(4x+4)} &&
\end{calculs}
\rem{On peut encore sortir des facteurs : $10x+6=2(5x+3)$ et $4x+4=4(x+1)$, donc
$A=8(5x+3)(x+1)$.}
\cas{B}
\begin{calculs}
B &= (7x-3)^{2}-(2x-1)^{2} &&\justif{$a=7x-3$ et $b=2x-1$}\\
B &= \big[(7x-3)+(2x-1)\big]\big[(7x-3)-(2x-1)\big] &&\\
B &= (7x-3+2x-1)(7x-3-2x+1) &&\justif{$-(2x-1)=-2x+1$}\\
B &= \res{(9x-4)(5x-2)} &&
\end{calculs}
\cas{C}
\begin{calculs}
C &= (3x-2)^{2}-49 &&\justif{$49=7^{2}$ : $a=3x-2$ et $b=7$}\\
C &= \big[(3x-2)+7\big]\big[(3x-2)-7\big] &&\\
C &= \res{(3x+5)(3x-9)} &&
\end{calculs}
\rem{Et $3x-9=3(x-3)$ : $C=3(3x+5)(x-3)$.}
\cas{D}
\begin{calculs}
D &= 81-(2x-3)^{2} &&\justif{$81=9^{2}$ : $a=9$ et $b=2x-3$}\\
D &= \big[9+(2x-3)\big]\big[9-(2x-3)\big] &&\\
D &= (9+2x-3)(9-2x+3) &&\\
D &= \res{(2x+6)(12-2x)} &&
\end{calculs}
\cas{E}
\begin{calculs}
E &= (x+5)^{2}-(x+2)^{2} &&\justif{$a=x+5$ et $b=x+2$}\\
E &= \big[(x+5)+(x+2)\big]\big[(x+5)-(x+2)\big] &&\\
E &= (2x+7)(x+5-x-2) &&\\
E &= (2x+7)\times 3 &&\justif{les $x$ s'éliminent dans le second crochet}\\
E &= \res{3(2x+7)} &&
\end{calculs}
\rem{Erreur fréquente : oublier les parenthèses dans le crochet $\big[(7x-3)-(2x-1)\big]$
et écrire $7x-3-2x-1$. L'énoncé demande ces parenthèses pour cette raison.}
\end{exo}

\begin{exo}[Le « 1 » caché et le carré]
\explique{Un facteur seul vaut $1\times$ ce facteur : quand on l'enlève, il reste $1$ dans
le crochet. Un carré $(\ldots)^{2}$ est un produit $(\ldots)(\ldots)$ : il contient le
facteur commun.}{propriété 5 du cours}
\cas{A}
\begin{calculs}
A &= (5x+4)(3x-2)+(5x+4) &&\justif{le second terme est $1\times(5x+4)$}\\
A &= (5x+4)(3x-2)+1\times(5x+4) &&\\
A &= (5x+4)\big[(3x-2)+1\big] &&\\
A &= \res{(5x+4)(3x-1)} &&
\end{calculs}
\cas{B}
\begin{calculs}
B &= 3(x-1)(x+2)-(x-1) &&\justif{facteur commun $(x-1)$}\\
B &= 3(x-1)(x+2)-1\times(x-1) &&\justif{le « 1 » caché}\\
B &= (x-1)\big[3(x+2)-1\big] &&\\
B &= (x-1)(3x+6-1) &&\\
B &= \res{(x-1)(3x+5)} &&
\end{calculs}
\cas{C}
\begin{calculs}
C &= (6r-4)(2r-3)+(6r-4)^{2} &&\justif{$(6r-4)^{2}=(6r-4)(6r-4)$}\\
C &= (6r-4)(2r-3)+(6r-4)(6r-4) &&\\
C &= (6r-4)\big[(2r-3)+(6r-4)\big] &&\\
C &= \res{(6r-4)(8r-7)} &&
\end{calculs}
\rem{Et $6r-4=2(3r-2)$ : $C=2(3r-2)(8r-7)$.}
\cas{D}
\begin{calculs}
D &= (4m-2)^{2}-(4m-2)(3m+1) &&\justif{$(4m-2)^{2}=(4m-2)(4m-2)$}\\
D &= (4m-2)\big[(4m-2)-(3m+1)\big] &&\\
D &= (4m-2)(4m-2-3m-1) &&\\
D &= \res{(4m-2)(m-3)} &&
\end{calculs}
\cas{E}
\explique{$x+2$ apparaît trois fois : dans le carré, dans $4(x+2)$, et seul, en fin
d'expression ($x+2=1\times(x+2)$).}{les deux derniers termes, $x$ et $2$, forment $x+2=1\times(x+2)$}
\begin{calculs}
E &= (x+2)^{2}+4(x+2)+x+2 &&\\
E &= (x+2)(x+2)+4(x+2)+1\times(x+2) &&\\
E &= (x+2)\big[(x+2)+4+1\big] &&\\
E &= \res{(x+2)(x+7)} &&
\end{calculs}
\rem{Erreur fréquente : écrire $(5x+4)\big[(3x-2)+0\big]$ dans $A$. Contrôle avec $x=0$ :
$4\times(-2)+4=-4$ et $4\times(-1)=-4$.}
\end{exo}

\begin{exo}[Faire apparaître le facteur commun]
\explique{Le facteur commun est caché : on le fait apparaître en réordonnant les termes,
en reconnaissant un opposé ($3-x=-(x-3)$) ou en factorisant un seul morceau avec une
identité.}{propriété 5 et méthode 7 du cours}
\cas{A}
\begin{calculs}
A &= (3x-4)(x+2)+6(-4+3x) &&\justif{$-4+3x=3x-4$ : même parenthèse, termes dans l'autre ordre}\\
A &= (3x-4)(x+2)+6(3x-4) &&\\
A &= (3x-4)\big[(x+2)+6\big] &&\\
A &= \res{(3x-4)(x+8)} &&
\end{calculs}
\cas{B}
\begin{calculs}
B &= (5-x)^{2}+(x-5) &&\justif{$5-x$ est l'opposé de $x-5$, et deux opposés ont le même carré}\\
B &= (x-5)^{2}+1\times(x-5) &&\justif{$(5-x)^{2}=(x-5)^{2}$}\\
B &= (x-5)\big[(x-5)+1\big] &&\\
B &= \res{(x-5)(x-4)} &&
\end{calculs}
\cas{C}
\begin{calculs}
C &= (-x+2)(3x+1)+(-3x-1)(x-7) &&\justif{$-3x-1=-(3x+1)$}\\
C &= (-x+2)(3x+1)-(3x+1)(x-7) &&\justif{le signe $-$ passe devant le produit}\\
C &= (3x+1)\big[(-x+2)-(x-7)\big] &&\\
C &= (3x+1)(-x+2-x+7) &&\\
C &= \res{(3x+1)(-2x+9)} &&
\end{calculs}
\cas{D}
\begin{calculs}
D &= 9x^{2}-16+(3x+4)(2x+5) &&\justif{$9x^{2}-16=(3x)^{2}-4^{2}$ fait apparaître $(3x+4)$}\\
D &= (3x+4)(3x-4)+(3x+4)(2x+5) &&\\
D &= (3x+4)\big[(3x-4)+(2x+5)\big] &&\\
D &= \res{(3x+4)(5x+1)} &&
\end{calculs}
\cas{E}
\begin{calculs}
E &= 16x^{2}+40x+25-2(4x+5) &&\justif{$16x^{2}+40x+25=(4x+5)^{2}$ : $a=4x$ et $b=5$ ; $2ab=40x$}\\
E &= (4x+5)^{2}-2(4x+5) &&\\
E &= (4x+5)\big[(4x+5)-2\big] &&\\
E &= \res{(4x+5)(4x+3)} &&
\end{calculs}
\rem{Erreur fréquente dans $C$ : écrire $(3x+1)\big[(-x+2)+(x-7)\big]$ en gardant le $+$.
Quand on remplace $(-3x-1)$ par $-(3x+1)$, le signe devant le produit devient $-$.}
\end{exo}

\begin{exo}[Choisir la méthode]
\explique{On applique la méthode 7 : facteur commun visible ? identité ? facteur commun
caché ?}{on n'oublie pas de réduire chaque facteur}
\cas{A}
\explique{$16x^{2}=(4x)^{2}$ et $49=7^{2}$ ; $2\times 4x\times 7=56x$ et le terme restant
est $-56x$. Donc $A=\res{(4x-7)^{2}}$.}{trois termes : carré parfait}
\cas{B}
\begin{calculs}
B &= 24x-10x^{2}+3x^{3} &&\justif{chaque terme contient $x$ ; les coefficients $24$ ; $10$ et $3$ n'ont pas de diviseur commun}\\
B &= x\times 24-x\times 10x+x\times 3x^{2} &&\\
B &= \res{x(24-10x+3x^{2})} &&
\end{calculs}
\cas{C}
\begin{calculs}
C &= (2x+5)(3x-4)+(2x+5)(4x+9) &&\justif{facteur commun $(2x+5)$}\\
C &= (2x+5)\big[(3x-4)+(4x+9)\big] &&\\
C &= \res{(2x+5)(7x+5)} &&
\end{calculs}
\cas{D}
\begin{calculs}
D &= (5x+3)(4x-9)-(5x+3)(6x-9) &&\justif{facteur commun $(5x+3)$}\\
D &= (5x+3)\big[(4x-9)-(6x-9)\big] &&\\
D &= (5x+3)(4x-9-6x+9) &&\\
D &= (5x+3)(-2x) &&\justif{les nombres s'éliminent}\\
D &= \res{-2x(5x+3)} &&
\end{calculs}
\cas{E}
\begin{calculs}
E &= (4x-1)^{2}-(3x+2)^{2} &&\justif{$a=4x-1$ et $b=3x+2$}\\
E &= \big[(4x-1)+(3x+2)\big]\big[(4x-1)-(3x+2)\big] &&\\
E &= (4x-1+3x+2)(4x-1-3x-2) &&\\
E &= \res{(7x+1)(x-3)} &&
\end{calculs}
\end{exo}

\partie{Fractions rationnelles}

\begin{exo}[Valeurs interdites]
\explique{Une valeur interdite annule le dénominateur : on cherche pour quelles valeurs de
$x$ le dénominateur vaut $0$. Un produit est nul quand l'un de ses facteurs est nul
(collège).}{définition 3 du cours}
\cas{A}
\explique{$x+3=0$ pour $x=-3$ : la valeur interdite est \res{$-3$}.}{contrôle :
$-3+3=0$}
\cas{B}
\explique{$3x-2=0$ donne $3x=2$, soit $x=\frac23$ : la valeur interdite est
\res{$\frac23$}.}{on ajoute $2$, puis on divise par $3$}
\cas{C}
\explique{$x(x+5)=0$ pour $x=0$ ou $x+5=0$, soit $x=-5$ : les valeurs interdites sont
\res{$0$ et $-5$}.}{un facteur nul suffit}
\cas{D}
\explique{$(3x+2)(2-x)=0$ pour $3x+2=0$, soit $x=-\frac23$, ou $2-x=0$, soit $x=2$ : les
valeurs interdites sont \res{$-\frac23$ et $2$}.}{une valeur par facteur}
\cas{E}
\explique{$x^{2}-16=(x+4)(x-4)$, nul pour $x=-4$ ou $x=4$ : les valeurs interdites sont
\res{$-4$ et $4$}.}{troisième identité ; $4^{2}$ et $(-4)^{2}$ valent $16$}
\rem{Le numérateur ne joue aucun rôle : seul le dénominateur peut rendre le calcul
impossible.}
\end{exo}

\begin{exo}[Premiers calculs]
\explique{On écrit chaque terme avec le dénominateur commun, puis on additionne les
numérateurs seulement.}{méthode 8 du cours ; ici $x\ne 0$}
\cas{A}
\begin{calculs}
A &= 2+\dfrac{1}{x} &&\justif{$2=\frac{2}{1}$ ; dénominateur commun $x$}\\
A &= \dfrac{2\fois{x}}{1\fois{x}}+\dfrac{1}{x} &&\\
A &= \res{\dfrac{2x+1}{x}} &&
\end{calculs}
\cas{B}
\begin{calculs}
B &= 4-\dfrac{3}{x} &&\\
B &= \dfrac{4\fois{x}}{1\fois{x}}-\dfrac{3}{x} &&\\
B &= \res{\dfrac{4x-3}{x}} &&
\end{calculs}
\cas{C}
\begin{calculs}
C &= \dfrac{2}{x}+\dfrac{5}{x} &&\justif{même dénominateur : on additionne les numérateurs}\\
C &= \res{\dfrac{7}{x}} &&
\end{calculs}
\cas{D}
\begin{calculs}
D &= x+\dfrac{3}{x} &&\\
D &= \dfrac{x\fois{x}}{1\fois{x}}+\dfrac{3}{x} &&\justif{$x\times x=x^{2}$}\\
D &= \res{\dfrac{x^{2}+3}{x}} &&
\end{calculs}
\cas{E}
\begin{calculs}
E &= \dfrac13+\dfrac{2}{x} &&\justif{dénominateur commun $3x$}\\
E &= \dfrac{1\fois{x}}{3\fois{x}}+\dfrac{2\fois{3}}{x\fois{3}} &&\\
E &= \res{\dfrac{x+6}{3x}} &&
\end{calculs}
\rem{Erreur fréquente : $2+\frac1x=\frac{3}{x}$. On n'ajoute pas un entier au numérateur :
$2$ s'écrit d'abord $\frac{2x}{x}$. Contrôle avec $x=1$ : $2+1=3$ et $\frac{2+1}{1}=3$.}
\end{exo}

\begin{exo}[Une seule fraction]
\explique{On donne d'abord les valeurs interdites (un dénominateur ne doit pas être nul),
on met au même dénominateur (le produit des dénominateurs, laissé factorisé), puis on
réduit le numérateur seulement.}{méthode 8 du cours}
\cas{A}
\explique{$x-2=0$ pour $x=2$ : on travaille pour $x\ne 2$.}{valeur interdite \res{$2$}}
\begin{calculs}
A &= 1+\dfrac{3}{x-2} &&\\
A &= \dfrac{x-2}{x-2}+\dfrac{3}{x-2} &&\justif{$1=\frac{x-2}{x-2}$}\\
A &= \dfrac{x-2+3}{x-2} &&\\
A &= \res{\dfrac{x+1}{x-2}} &&
\end{calculs}
\cas{B}
\explique{Les dénominateurs $x$ et $x+2$ s'annulent en $0$ et en $-2$ : on travaille pour
$x\ne 0$ et $x\ne -2$.}{valeurs interdites \res{$0$ et $-2$}}
\begin{calculs}
B &= \dfrac{4}{x}-\dfrac{3}{x+2} &&\justif{dénominateur commun $x(x+2)$}\\
B &= \dfrac{4\fois{(x+2)}}{x\fois{(x+2)}}-\dfrac{3\fois{x}}{(x+2)\fois{x}} &&\\
B &= \dfrac{4x+8-3x}{x(x+2)} &&\\
B &= \res{\dfrac{x+8}{x(x+2)}} &&
\end{calculs}
\cas{C}
\explique{$x+3=0$ pour $x=-3$, et $x=0$ annule le second dénominateur : on travaille pour
$x\ne -3$ et $x\ne 0$.}{valeurs interdites \res{$-3$ et $0$}}
\begin{calculs}
C &= \dfrac{x}{x+3}+\dfrac{2}{x} &&\justif{dénominateur commun $x(x+3)$}\\
C &= \dfrac{x\fois{x}}{(x+3)\fois{x}}+\dfrac{2\fois{(x+3)}}{x\fois{(x+3)}} &&\\
C &= \res{\dfrac{x^{2}+2x+6}{x(x+3)}} &&\justif{$x^{2}+2(x+3)=x^{2}+2x+6$}
\end{calculs}
\cas{D}
\explique{$x-1=0$ pour $x=1$ : on travaille pour $x\ne 1$.}{valeur interdite \res{$1$}}
\begin{calculs}
D &= 3-\dfrac{x+2}{x-1} &&\\
D &= \dfrac{3\fois{(x-1)}}{1\fois{(x-1)}}-\dfrac{x+2}{x-1} &&\\
D &= \dfrac{3(x-1)-(x+2)}{x-1} &&\justif{le $-$ porte sur tout le numérateur $x+2$}\\
D &= \dfrac{3x-3-x-2}{x-1} &&\\
D &= \res{\dfrac{2x-5}{x-1}} &&
\end{calculs}
\rem{Erreur fréquente dans $D$ : écrire $3x-3-x+2$, sans parenthèses autour de $x+2$.}
\cas{E}
\explique{$x-2$ et $x+2$ s'annulent en $2$ et en $-2$ : on travaille pour $x\ne 2$ et
$x\ne -2$.}{valeurs interdites \res{$2$ et $-2$}}
\begin{calculs}
E &= \dfrac{1}{x-2}+\dfrac{1}{x+2} &&\justif{dénominateur commun $(x-2)(x+2)$}\\
E &= \dfrac{1\fois{(x+2)}}{(x-2)\fois{(x+2)}}+\dfrac{1\fois{(x-2)}}{(x+2)\fois{(x-2)}} &&\\
E &= \res{\dfrac{2x}{(x-2)(x+2)}} &&\justif{$x+2+x-2=2x$}
\end{calculs}
\rem{On ne simplifie que des facteurs communs au numérateur et au dénominateur : dans
$\frac{x+1}{x-2}$, les $x$ ne se simplifient pas.}
\end{exo}

\begin{exo}[Démontrer une égalité]
\explique{Pour démontrer une égalité, on part d'un membre (le plus compliqué, ou celui
qui a plusieurs fractions) et on le transforme jusqu'à obtenir l'autre.}{on ne part pas
de l'égalité à démontrer}
\cas{1)}
\explique{Pour $x\ne 1$, on transforme le membre de droite.}{il contient deux termes à
réunir}
\begin{calculs}
2-\dfrac{1}{x-1} &= \dfrac{2\fois{(x-1)}}{1\fois{(x-1)}}-\dfrac{1}{x-1} &&\justif{$2=\frac{2}{1}$}\\
2-\dfrac{1}{x-1} &= \dfrac{2x-2-1}{x-1} &&\\
2-\dfrac{1}{x-1} &= \dfrac{2x-3}{x-1} &&
\end{calculs}
\explique{On obtient le membre de gauche : pour tout $x\ne 1$,
$\res{\dfrac{2x-3}{x-1}=2-\dfrac{1}{x-1}}$.}{on conclut par une phrase}
\cas{2)}
\explique{Pour $x\ne 0$ et $x\ne 2$, on transforme le membre de gauche.}{dénominateur
commun $x(x-2)$}
\begin{calculs}
\dfrac{1}{x-2}-\dfrac{1}{x} &= \dfrac{1\fois{x}}{(x-2)\fois{x}}-\dfrac{1\fois{(x-2)}}{x\fois{(x-2)}} &&\\
\dfrac{1}{x-2}-\dfrac{1}{x} &= \dfrac{x-(x-2)}{x(x-2)} &&\justif{parenthèses : le $-$ porte sur $x-2$}\\
\dfrac{1}{x-2}-\dfrac{1}{x} &= \dfrac{x-x+2}{x(x-2)} &&\\
\dfrac{1}{x-2}-\dfrac{1}{x} &= \dfrac{2}{x(x-2)} &&
\end{calculs}
\explique{Pour tout $x\ne 0$ et $x\ne 2$,
$\res{\dfrac{1}{x-2}-\dfrac{1}{x}=\dfrac{2}{x(x-2)}}$.}{}
\end{exo}

\partie{Équations du premier degré}

\begin{exo}[Les deux règles]
\explique{On regroupe les $x$ d'un côté et les nombres de l'autre en ajoutant ou en
soustrayant le même terme aux deux membres, puis on divise par le coefficient de
$x$.}{propriété 7 et méthode 9 du cours}
\cas{a)}
\begin{calculs}
4x+3 &= 3x+8 &&\justif{on recopie}\\
4x-3x &= 8-3 &&\justif{on soustrait $3x$ et $3$ aux deux membres}\\
x &= 5 &&
\end{calculs}
\explique{$\res{S=\ens{5}}$.}{vérification : $4\times 5+3=23$ et $3\times 5+8=23$}
\cas{b)}
\begin{calculs}
3x+2 &= 4x-6 &&\\
2+6 &= 4x-3x &&\justif{les $x$ du côté où il y en a le plus : on évite les signes moins}\\
8 &= x &&
\end{calculs}
\explique{$\res{S=\ens{8}}$.}{vérification : $26=26$}
\cas{c)}
\begin{calculs}
2x+5 &= 6x+1 &&\\
5-1 &= 6x-2x &&\\
4 &= 4x &&\\
x &= 1 &&\justif{on divise par $4$}
\end{calculs}
\explique{$\res{S=\ens{1}}$.}{vérification : $7=7$}
\cas{d)}
\begin{calculs}
7-2x &= 0 &&\\
7 &= 2x &&\justif{on ajoute $2x$ aux deux membres}\\
x &= \dfrac72 &&\justif{on divise par $2$}
\end{calculs}
\explique{$\res{S=\ens{\frac72}}$.}{on garde la fraction (ou $3{,}5$)}
\cas{e)}
\begin{calculs}
2x-9 &= 5x+12 &&\\
-9-12 &= 5x-2x &&\\
-21 &= 3x &&\\
x &= -7 &&\justif{on divise par $3$ : $-21\div 3=-7$}
\end{calculs}
\explique{$\res{S=\ens{-7}}$.}{vérification : $-14-9=-23$ et $-35+12=-23$}
\cas{f)}
\begin{calculs}
3x+7 &= 3x-4 &&\\
3x-3x &= -4-7 &&\\
0x &= -11 &&\justif{les $x$ disparaissent}
\end{calculs}
\explique{$0x=-11$ n'est vrai pour aucun $x$ : $\res{S=\varnothing}$.}{$0\times x$ vaut
toujours $0$, jamais $-11$}
\rem{Erreur fréquente : diviser par $3$ avant d'avoir regroupé, ou changer le signe d'un
terme qui ne change pas de membre.}
\end{exo}

\begin{exo}[Avec des parenthèses]
\explique{On développe d'abord chaque membre (un $-$ devant une parenthèse change tous les
signes), on réduit, puis on applique les deux règles.}{méthode 9 du cours}
\cas{a)}
\begin{calculs}
7-(x-2) &= 5x-(4x-9) &&\justif{on recopie}\\
7-x+2 &= 5x-4x+9 &&\justif{les deux $-$ changent les signes des parenthèses}\\
9-x &= x+9 &&\justif{on réduit chaque membre}\\
9-9 &= x+x &&\\
0 &= 2x &&\\
x &= 0 &&
\end{calculs}
\explique{$\res{S=\ens{0}}$.}{vérification : $7+2=9$ et $0+9=9$}
\cas{b)}
\begin{calculs}
4(x-2)+3(3-x) &= 0 &&\\
4x-8+9-3x &= 0 &&\\
x+1 &= 0 &&\\
x &= -1 &&
\end{calculs}
\explique{$\res{S=\ens{-1}}$.}{vérification : $4\times(-3)+3\times 4=0$}
\cas{c)}
\begin{calculs}
6(x+3)-2(x+4) &= x+16 &&\\
6x+18-2x-8 &= x+16 &&\justif{$-2\times 4=-8$}\\
4x+10 &= x+16 &&\\
4x-x &= 16-10 &&\\
3x &= 6 &&\\
x &= 2 &&
\end{calculs}
\explique{$\res{S=\ens{2}}$.}{vérification : $30-12=18$ et $2+16=18$}
\cas{d)}
\begin{calculs}
3(x-2)-2(x+1) &= 5(x-4) &&\\
3x-6-2x-2 &= 5x-20 &&\\
x-8 &= 5x-20 &&\\
-8+20 &= 5x-x &&\\
12 &= 4x &&\\
x &= 3 &&
\end{calculs}
\explique{$\res{S=\ens{3}}$.}{vérification : $3-8=-5$ et $5\times(-1)=-5$}
\cas{e)}
\begin{calculs}
m-(4-2m) &= 2m+5 &&\justif{l'inconnue s'appelle $m$}\\
m-4+2m &= 2m+5 &&\\
3m-4 &= 2m+5 &&\\
3m-2m &= 5+4 &&\\
m &= 9 &&
\end{calculs}
\explique{$\res{S=\ens{9}}$.}{vérification : $9-(4-18)=23$ et $18+5=23$}
\end{exo}

\begin{exo}[Coefficients fractionnaires]
\explique{L'équation a la forme $ax=b$ avec $a$ fractionnaire : on divise par $a$, c'est-à-dire
qu'on multiplie par son inverse (chapitre 1).}{l'inverse de $\frac{p}{q}$ est $\frac{q}{p}$}
\cas{a)}
\begin{calculs}
\dfrac{-5x}{3} &= \dfrac{10}{3} &&\justif{même dénominateur}\\
-5x &= 10 &&\justif{on multiplie les deux membres par $3$}\\
x &= -2 &&\justif{on divise par $-5$}
\end{calculs}
\explique{$\res{S=\ens{-2}}$.}{}
\cas{b)}
\begin{calculs}
\dfrac{12x}{5} &= \dfrac{3}{10} &&\\
x &= \dfrac{3}{10}\times\dfrac{5}{12} &&\justif{on multiplie par l'inverse de $\frac{12}{5}$}\\
x &= \dfrac{\barre{3}\times\barre{5}}{2\times\barre{5}\times 4\times\barre{3}} &&\justif{$10=2\times 5$ et $12=4\times 3$}\\
x &= \dfrac18 &&
\end{calculs}
\explique{$\res{S=\ens{\frac18}}$.}{}
\cas{c)}
\begin{calculs}
\dfrac34x &= \dfrac92 &&\\
x &= \dfrac92\times\dfrac43 &&\justif{l'inverse de $\frac34$ est $\frac43$}\\
x &= \dfrac{3\times\barre{3}\times 2\times\barre{2}}{\barre{2}\times\barre{3}} &&\justif{$9=3\times 3$ et $4=2\times 2$}\\
x &= 6 &&
\end{calculs}
\explique{$\res{S=\ens{6}}$.}{vérification : $\frac34\times 6=\frac{18}{4}=\frac92$}
\cas{d)}
\begin{calculs}
-3x &= \dfrac{12}{5} &&\\
x &= \dfrac{12}{5}\times\left(-\dfrac13\right) &&\justif{diviser par $-3$, c'est multiplier par $-\frac13$}\\
x &= -\dfrac{4\times\barre{3}}{5\times\barre{3}} &&\\
x &= -\dfrac45 &&
\end{calculs}
\explique{$\res{S=\ens{-\frac45}}$.}{}
\cas{e)}
\begin{calculs}
\dfrac{-3}{7}x &= \dfrac{-1}{14} &&\\
x &= \dfrac{-1}{14}\times\dfrac{7}{-3} &&\justif{deux signes moins : le résultat est positif}\\
x &= \dfrac{1\times\barre{7}}{2\times\barre{7}\times 3} &&\justif{$14=2\times 7$}\\
x &= \dfrac16 &&
\end{calculs}
\explique{$\res{S=\ens{\frac16}}$.}{}
\cas{f)}
\begin{calculs}
\dfrac56 &= \dfrac{-10}{9}x &&\justif{l'inconnue est à droite}\\
x &= \dfrac56\times\dfrac{9}{-10} &&\\
x &= -\dfrac{\barre{5}\times 3\times\barre{3}}{2\times\barre{3}\times 2\times\barre{5}} &&\justif{$6=2\times 3$ ; $9=3\times 3$ ; $10=2\times 5$}\\
x &= -\dfrac34 &&
\end{calculs}
\explique{$\res{S=\ens{-\frac34}}$.}{}
\rem{Erreur fréquente : soustraire le coefficient au lieu de diviser, par exemple écrire
$x=\frac92-\frac34$ en c). Le coefficient multiplie $x$ : on divise.}
\end{exo}

\begin{exo}[Supprimer les fractions]
\explique{On multiplie \emph{chaque terme} des deux membres par le dénominateur commun :
les fractions disparaissent. Un numérateur qui est une somme se met entre
parenthèses.}{méthode 9 du cours, exemple b)}
\cas{a)}
\begin{calculs}
-\dfrac13x+4 &= x-8 &&\justif{dénominateur $3$}\\
3\times\left(-\dfrac13x\right)+3\times 4 &= 3\times x-3\times 8 &&\justif{chaque terme $\times 3$}\\
-x+12 &= 3x-24 &&\\
12+24 &= 3x+x &&\\
36 &= 4x &&\\
x &= 9 &&
\end{calculs}
\explique{$\res{S=\ens{9}}$.}{vérification : $-3+4=1$ et $9-8=1$}
\cas{b)}
\begin{calculs}
\dfrac52x+1 &= 3x-4 &&\justif{chaque terme $\times 2$}\\
5x+2 &= 6x-8 &&\\
2+8 &= 6x-5x &&\\
x &= 10 &&
\end{calculs}
\explique{$\res{S=\ens{10}}$.}{vérification : $26=26$}
\cas{c)}
\begin{calculs}
\dfrac{3x}{5}-\dfrac43 &= \dfrac{7}{15} &&\justif{$15$ est un multiple de $5$ et de $3$}\\
15\times\dfrac{3x}{5}-15\times\dfrac43 &= 15\times\dfrac{7}{15} &&\\
9x-20 &= 7 &&\justif{$15\div 5=3$ et $3\times 3x=9x$ ; $15\div 3=5$ et $5\times 4=20$}\\
9x &= 27 &&\\
x &= 3 &&
\end{calculs}
\explique{$\res{S=\ens{3}}$.}{vérification : $\frac95-\frac43=\frac{27-20}{15}=\frac{7}{15}$}
\cas{d)}
\begin{calculs}
\dfrac{x}{2}+\dfrac53 &= -\dfrac{x}{6}+9 &&\justif{dénominateur commun $6$ ; le $9$ aussi est multiplié}\\
3x+10 &= -x+54 &&\\
3x+x &= 54-10 &&\\
4x &= 44 &&\\
x &= 11 &&
\end{calculs}
\explique{$\res{S=\ens{11}}$.}{vérification : $\frac{11}{2}+\frac53=\frac{43}{6}$ et $-\frac{11}{6}+\frac{54}{6}=\frac{43}{6}$}
\cas{e)}
\begin{calculs}
\dfrac{x-3}{2}-4 &= \dfrac{3x+1}{4}+\dfrac12 &&\justif{dénominateur commun $4$}\\
2(x-3)-16 &= 3x+1+2 &&\justif{$4\times\frac{x-3}{2}=2(x-3)$ ; $4\times 4=16$ ; $4\times\frac12=2$}\\
2x-6-16 &= 3x+3 &&\\
2x-22 &= 3x+3 &&\\
-22-3 &= 3x-2x &&\\
x &= -25 &&
\end{calculs}
\explique{$\res{S=\ens{-25}}$.}{vérification : les deux membres valent $-18$}
\cas{f)}
\begin{calculs}
\dfrac{3x+1}{4}-\dfrac{x-3}{4} &= \dfrac{x+5}{3} &&\justif{dénominateur commun $12$}\\
3(3x+1)-3(x-3) &= 4(x+5) &&\justif{parenthèses autour de chaque numérateur}\\
9x+3-3x+9 &= 4x+20 &&\justif{$-3\times(-3)=+9$}\\
6x+12 &= 4x+20 &&\\
6x-4x &= 20-12 &&\\
2x &= 8 &&\\
x &= 4 &&
\end{calculs}
\explique{$\res{S=\ens{4}}$.}{vérification : $\frac{13-1}{4}=3$ et $\frac93=3$}
\rem{Erreur fréquente : oublier de multiplier le terme sans fraction (le $-4$ de e)
devient $-16$), ou oublier la parenthèse : $3(x-3)$, et non $3x-3$.}
\end{exo}

\begin{exo}[Produit en croix, cas particuliers]
\cas{1) a)}
\begin{calculs}
\dfrac{3x+2}{4} &= \dfrac{5x-1}{3} &&\justif{deux fractions égales}\\
3(3x+2) &= 4(5x-1) &&\justif{produit en croix}\\
9x+6 &= 20x-4 &&\\
6+4 &= 20x-9x &&\\
10 &= 11x &&\\
x &= \dfrac{10}{11} &&
\end{calculs}
\explique{$\res{S=\ens{\frac{10}{11}}}$.}{on garde la valeur exacte}
\cas{1) b)}
\begin{calculs}
\dfrac{3x-2}{5} &= \dfrac23 &&\\
3(3x-2) &= 5\times 2 &&\justif{produit en croix}\\
9x-6 &= 10 &&\\
9x &= 16 &&\\
x &= \dfrac{16}{9} &&
\end{calculs}
\explique{$\res{S=\ens{\frac{16}{9}}}$.}{}
\cas{2) a)}
\begin{calculs}
\dfrac12(x+4)-\dfrac23(x-3) &= \dfrac16(-x+12)+2 &&\justif{chaque terme $\times 6$}\\
3(x+4)-4(x-3) &= -x+12+12 &&\justif{$6\times\frac12=3$ ; $6\times\frac23=4$ ; $6\times 2=12$}\\
3x+12-4x+12 &= -x+24 &&\\
-x+24 &= -x+24 &&\justif{les deux membres sont identiques}\\
0x &= 0 &&
\end{calculs}
\explique{$0x=0$ est vrai pour tout $x$ : $\res{S=\R}$.}{propriété 7 : $a=0$ et $b=0$}
\cas{2) b)}
\begin{calculs}
\dfrac{x+1}{3}-\dfrac{2x-1}{4}-1 &= -\dfrac{2x+1}{12} &&\justif{chaque terme $\times 12$}\\
4(x+1)-3(2x-1)-12 &= -(2x+1) &&\\
4x+4-6x+3-12 &= -2x-1 &&\\
-2x-5 &= -2x-1 &&\\
0x &= 4 &&\justif{on ajoute $2x$ et $5$}
\end{calculs}
\explique{$0x=4$ n'est jamais vrai : $\res{S=\varnothing}$.}{propriété 7 : $a=0$ et
$b\ne 0$}
\rem{Quand les $x$ disparaissent, on ne s'arrête pas : on lit l'égalité qui reste. Vraie
($0=0$), tous les réels sont solutions ; fausse ($0=4$), il n'y en a aucune.}
\end{exo}

\begin{exo}[L'inconnue au dénominateur]
\explique{On donne la valeur interdite (dénominateur nul), on fait un produit en croix
($2=\frac21$), on résout, puis on vérifie que la solution n'est pas
interdite.}{méthode 10 du cours}
\cas{a)}
\explique{$2x+1=0$ pour $x=-\frac12$. Pour $x\ne -\frac12$ :}{on écarte $-\frac12$}
\begin{calculs}
\dfrac{3x-2}{2x+1} &= 2 &&\\
3x-2 &= 2(2x+1) &&\justif{produit en croix avec $\frac21$}\\
3x-2 &= 4x+2 &&\\
-2-2 &= 4x-3x &&\\
x &= -4 &&
\end{calculs}
\explique{$-4\ne -\frac12$ : $\res{S=\ens{-4}}$.}{vérification : $\frac{-14}{-7}=2$}
\cas{b)}
\explique{$2-x=0$ pour $x=2$. Pour $x\ne 2$ :}{}
\begin{calculs}
\dfrac{3-2x}{2-x} &= 1 &&\\
3-2x &= 2-x &&\justif{produit en croix avec $\frac11$}\\
3-2 &= 2x-x &&\\
x &= 1 &&
\end{calculs}
\explique{$1\ne 2$ : $\res{S=\ens{1}}$.}{vérification : $\frac11=1$}
\cas{c)}
\explique{$2x+1=0$ pour $x=-\frac12$. Pour $x\ne -\frac12$ :}{}
\begin{calculs}
\dfrac{3-x}{2x+1} &= \dfrac34 &&\\
4(3-x) &= 3(2x+1) &&\justif{produit en croix}\\
12-4x &= 6x+3 &&\\
12-3 &= 6x+4x &&\\
9 &= 10x &&\\
x &= \dfrac{9}{10} &&
\end{calculs}
\explique{$\frac{9}{10}$ n'est pas interdit : $\res{S=\ens{\frac{9}{10}}}$.}{}
\cas{d)}
\explique{$1+2x=0$ pour $x=-\frac12$. Pour $x\ne -\frac12$ :}{}
\begin{calculs}
\dfrac{2+x}{1+2x} &= \dfrac12 &&\\
2(2+x) &= 1+2x &&\justif{produit en croix}\\
4+2x &= 1+2x &&\\
0x &= -3 &&\justif{les $x$ disparaissent}
\end{calculs}
\explique{$0x=-3$ est toujours faux : $\res{S=\varnothing}$.}{}
\cas{e)}
\explique{$x-2=0$ pour $x=2$. Pour $x\ne 2$ :}{}
\begin{calculs}
\dfrac{3-x}{x-2} &= 2 &&\\
3-x &= 2(x-2) &&\\
3-x &= 2x-4 &&\\
3+4 &= 2x+x &&\\
7 &= 3x &&\\
x &= \dfrac73 &&
\end{calculs}
\explique{$\frac73\ne 2$ : $\res{S=\ens{\frac73}}$.}{vérification : $\frac{2/3}{1/3}=2$}
\cas{f)}
\explique{$2+x=0$ pour $x=-2$. Pour $x\ne -2$ :}{}
\begin{calculs}
\dfrac{-3x}{2+x} &= 2 &&\\
-3x &= 2(2+x) &&\\
-3x &= 4+2x &&\\
-3x-2x &= 4 &&\\
-5x &= 4 &&\\
x &= -\dfrac45 &&
\end{calculs}
\explique{$-\frac45\ne -2$ : $\res{S=\ens{-\frac45}}$.}{}
\rem{Le dernier contrôle n'est pas une formalité : une solution qui annule le
dénominateur serait à rejeter.}
\end{exo}

\partie{Équations produit}

\begin{exo}[Produit nul]
\explique{Un produit est nul si et seulement si l'un au moins de ses facteurs est nul : on
écrit \og facteur $=0$\fg{} pour chaque facteur qui contient $x$.}{propriété 8 du cours}
\cas{a)}
\explique{$2x(3x-5)=0$ équivaut à $2x=0$ ou $3x-5=0$, soit $x=0$ ou $x=\frac53$ :
$\res{S=\ens{0\,;\frac53}}$.}{deux facteurs : deux équations du premier degré}
\cas{b)}
\explique{$-4x^{2}(3x-8)=0$ équivaut à $-4x^{2}=0$ ou $3x-8=0$. Or $-4x^{2}=0$ seulement
pour $x=0$ ; et $3x-8=0$ donne $x=\frac83$ : $\res{S=\ens{0\,;\frac83}}$.}{$-4$ n'est
jamais nul ; $x^{2}=x\times x$ est nul pour $x=0$}
\cas{c)}
\explique{$(-2x+5)^{2}=(-2x+5)(-2x+5)$ : les deux facteurs sont les mêmes. $-2x+5=0$ donne
$x=\frac52$ : $\res{S=\ens{\frac52}}$.}{une seule solution}
\cas{d)}
\explique{$x^{2}(x+2)(-4-3x)=0$ équivaut à $x^{2}=0$ ou $x+2=0$ ou $-4-3x=0$, soit $x=0$
ou $x=-2$ ou $x=-\frac43$.}{$-4-3x=0$ donne $-3x=4$}
\explique{$\res{S=\ens{-2\,;-\frac43\,;0}}$.}{on range les solutions dans l'ordre croissant}
\cas{e)}
\explique{$x+1=0$ ou $3x-2=0$ ou $x-5=0$, soit $x=-1$ ou $x=\frac23$ ou $x=5$ :
$\res{S=\ens{-1\,;\frac23\,;5}}$.}{trois facteurs : trois solutions}
\rem{Erreur fréquente : développer un produit déjà factorisé. Avec un produit nul, la forme
factorisée donne directement les solutions.}
\end{exo}

\begin{exo}[Factoriser d'abord]
\explique{On met tout dans le membre de gauche ($\ldots=0$), on factorise, puis on applique
le produit nul.}{méthode 11 du cours}
\cas{a)}
\begin{calculs}
3x^{2}-8x &= 0 &&\justif{$x$ est commun}\\
x(3x-8) &= 0 &&
\end{calculs}
\explique{$x=0$ ou $3x-8=0$ : $\res{S=\ens{0\,;\frac83}}$.}{}
\cas{b)}
\begin{calculs}
2x^{3}-5x^{2} &= 0 &&\justif{$x^{2}$ est commun : $x^{3}=x^{2}\times x$}\\
x^{2}(2x-5) &= 0 &&
\end{calculs}
\explique{$x^{2}=0$ ou $2x-5=0$ : $\res{S=\ens{0\,;\frac52}}$.}{}
\cas{c)}
\begin{calculs}
(x+3)^{2} &= (x+3)(4x-3) &&\justif{on recopie}\\
(x+3)^{2}-(x+3)(4x-3) &= 0 &&\justif{tout dans le membre de gauche}\\
(x+3)\big[(x+3)-(4x-3)\big] &= 0 &&\justif{facteur commun $(x+3)$}\\
(x+3)(-3x+6) &= 0 &&\justif{$x+3-4x+3=-3x+6$}
\end{calculs}
\explique{$x+3=0$ ou $-3x+6=0$, soit $x=-3$ ou $x=2$ : $\res{S=\ens{-3\,;2}}$.}{}
\rem{Ne jamais diviser par $(x+3)$ : on obtiendrait $x+3=4x-3$, soit seulement $x=2$, et
on perdrait la solution $-3$.}
\cas{d)}
\begin{calculs}
(2x-5)(3x+1)-(2x-5)(4-x) &= 0 &&\justif{tout à gauche}\\
(2x-5)\big[(3x+1)-(4-x)\big] &= 0 &&\\
(2x-5)(4x-3) &= 0 &&\justif{$3x+1-4+x=4x-3$}
\end{calculs}
\explique{$2x-5=0$ ou $4x-3=0$ : $\res{S=\ens{\frac34\,;\frac52}}$.}{}
\cas{e)}
\explique{Pas de facteur commun visible : on factorise le morceau $9x^{2}-4$ avec la
troisième identité, ce qui fait apparaître $(3x+2)$ et $(3x-2)$.}{propriété 5 : une
identité dans un seul morceau}
\begin{calculs}
9x^{2}-4-3(3x+2)+x(3x+2) &= 0 &&\\
(3x+2)(3x-2)-3(3x+2)+x(3x+2) &= 0 &&\justif{$9x^{2}-4=(3x)^{2}-2^{2}$}\\
(3x+2)\big[(3x-2)-3+x\big] &= 0 &&\justif{facteur commun $(3x+2)$}\\
(3x+2)(4x-5) &= 0 &&
\end{calculs}
\explique{$3x+2=0$ ou $4x-5=0$ : $\res{S=\ens{-\frac23\,;\frac54}}$.}{}
\cas{f)}
\begin{calculs}
(3x-1)(x+2)^{2}-9(3x-1) &= 0 &&\justif{tout à gauche}\\
(3x-1)\big[(x+2)^{2}-9\big] &= 0 &&\justif{facteur commun $(3x-1)$}\\
(3x-1)\big[(x+2)-3\big]\big[(x+2)+3\big] &= 0 &&\justif{$a^{2}-b^{2}$ avec $a=x+2$ et $b=3$}\\
(3x-1)(x-1)(x+5) &= 0 &&
\end{calculs}
\explique{$3x-1=0$ ou $x-1=0$ ou $x+5=0$ : $\res{S=\ens{-5\,;\frac13\,;1}}$.}{trois
facteurs}
\end{exo}

\begin{exo}[Différence de deux carrés nulle]
\explique{Le membre de gauche est de la forme $a^{2}-b^{2}$ : on le factorise en
$(a-b)(a+b)$, puis on applique le produit nul.}{méthode 11 du cours, exemple d)}
\cas{a)}
\begin{calculs}
9x^{2}-(x+2)^{2} &= 0 &&\justif{$9x^{2}=(3x)^{2}$ : $a=3x$ et $b=x+2$}\\
\big[3x-(x+2)\big]\big[3x+(x+2)\big] &= 0 &&\\
(2x-2)(4x+2) &= 0 &&
\end{calculs}
\explique{$2x-2=0$ ou $4x+2=0$, soit $x=1$ ou $x=-\frac12$ :
$\res{S=\ens{-\frac12\,;1}}$.}{}
\cas{b)}
\begin{calculs}
(3x+1)^{2}-(x+3)^{2} &= 0 &&\justif{$a=3x+1$ et $b=x+3$}\\
\big[(3x+1)-(x+3)\big]\big[(3x+1)+(x+3)\big] &= 0 &&\\
(2x-2)(4x+4) &= 0 &&
\end{calculs}
\explique{$2x-2=0$ ou $4x+4=0$ : $\res{S=\ens{-1\,;1}}$.}{}
\cas{c)}
\begin{calculs}
(3x+5)^{2}-4(x+1)^{2} &= 0 &&\justif{$4(x+1)^{2}=\big[2(x+1)\big]^{2}$ : $b=2(x+1)$}\\
\big[(3x+5)-2(x+1)\big]\big[(3x+5)+2(x+1)\big] &= 0 &&\\
(3x+5-2x-2)(3x+5+2x+2) &= 0 &&\\
(x+3)(5x+7) &= 0 &&
\end{calculs}
\explique{$x+3=0$ ou $5x+7=0$ : $\res{S=\ens{-3\,;-\frac75}}$.}{}
\cas{d)}
\begin{calculs}
9(x+1)^{2}-4(2x+3)^{2} &= 0 &&\justif{$a=3(x+1)$ et $b=2(2x+3)$}\\
\big[3(x+1)-2(2x+3)\big]\big[3(x+1)+2(2x+3)\big] &= 0 &&\\
(3x+3-4x-6)(3x+3+4x+6) &= 0 &&\\
(-x-3)(7x+9) &= 0 &&
\end{calculs}
\explique{$-x-3=0$ ou $7x+9=0$ : $\res{S=\ens{-3\,;-\frac97}}$.}{}
\rem{Erreur fréquente : $4(x+1)^{2}=(4x+4)^{2}$. Le $4$ est le carré de $2$ : c'est
$\big(2(x+1)\big)^{2}=(2x+2)^{2}$.}
\end{exo}

\begin{exo}[Développer d'abord]
\explique{Les deux membres contiennent des $x^{2}$ (ou des $x^{3}$) qui vont s'éliminer :
on développe, on regroupe, et il reste une équation du premier degré, ou une équation
produit.}{méthode 9 du cours}
\cas{a)}
\begin{calculs}
(3x-2)(4+x) &= 3x^{2}+5 &&\\
12x+3x^{2}-8-2x &= 3x^{2}+5 &&\justif{double distributivité}\\
3x^{2}+10x-8 &= 3x^{2}+5 &&\\
10x &= 13 &&\justif{les $3x^{2}$ s'éliminent ; on ajoute $8$}\\
x &= \dfrac{13}{10} &&
\end{calculs}
\explique{$\res{S=\ens{\frac{13}{10}}}$.}{soit $1{,}3$}
\cas{b)}
\begin{calculs}
(2x+3)^{2} &= 4x(x+0{,}5) &&\\
4x^{2}+12x+9 &= 4x^{2}+2x &&\justif{$(a+b)^{2}$ à gauche ; $4x\times 0{,}5=2x$}\\
12x-2x &= -9 &&\justif{les $4x^{2}$ s'éliminent}\\
10x &= -9 &&\\
x &= -\dfrac{9}{10} &&
\end{calculs}
\explique{$\res{S=\ens{-\frac{9}{10}}}$.}{soit $-0{,}9$}
\cas{c)}
\begin{calculs}
(x+3)(x-1) &= (x+5)(x-4) &&\\
x^{2}-x+3x-3 &= x^{2}-4x+5x-20 &&\\
x^{2}+2x-3 &= x^{2}+x-20 &&\\
2x-x &= -20+3 &&\justif{les $x^{2}$ s'éliminent}\\
x &= -17 &&
\end{calculs}
\explique{$\res{S=\ens{-17}}$.}{vérification : $(-14)\times(-18)=252$ et
$(-12)\times(-21)=252$}
\cas{d)}
\begin{calculs}
x(4-x)^{2} &= x^{2}(x+3) &&\\
x(16-8x+x^{2}) &= x^{3}+3x^{2} &&\justif{$(a-b)^{2}$ avec $a=4$ et $b=x$}\\
16x-8x^{2}+x^{3} &= x^{3}+3x^{2} &&\\
16x-8x^{2}-3x^{2} &= 0 &&\justif{les $x^{3}$ s'éliminent ; tout à gauche}\\
16x-11x^{2} &= 0 &&\\
x(16-11x) &= 0 &&\justif{$x$ est commun}
\end{calculs}
\explique{$x=0$ ou $16-11x=0$ : $\res{S=\ens{0\,;\frac{16}{11}}}$.}{ici, il reste une
équation produit}
\rem{En d), on ne divise pas par $x$ : on perdrait la solution $0$.}
\end{exo}

\begin{exo}[Choisir la bonne écriture]
\cas{1)}
\begin{calculs}
E(x) &= (x-2)^{2}-36 &&\justif{forme A}\\
E(x) &= x^{2}-4x+4-36 &&\justif{$(a-b)^{2}$ avec $a=x$ et $b=2$}\\
E(x) &= x^{2}-4x-32 &&\justif{forme B}
\end{calculs}
\begin{calculs}
E(x) &= (x-2)^{2}-6^{2} &&\justif{$a^{2}-b^{2}$ avec $a=x-2$ et $b=6$}\\
E(x) &= \big[(x-2)-6\big]\big[(x-2)+6\big] &&\\
E(x) &= (x-8)(x+4) &&\justif{forme C}
\end{calculs}
\explique{On a bien $\res{E(x)=x^{2}-4x-32}$ et $\res{E(x)=(x-8)(x+4)}$.}{contrôle :
$(x-8)(x+4)=x^{2}+4x-8x-32$}
\cas{2)}
\explique{$E(0)$ : avec la forme B, $E(0)=0-0-32=\res{-32}$.}{$x=0$ annule les termes en
$x$}
\explique{$E(x)=0$ : avec la forme C, $(x-8)(x+4)=0$, soit $x=8$ ou $x=-4$ :
$\res{S=\ens{-4\,;8}}$.}{forme factorisée et produit nul}
\explique{$E(x)=-32$ : avec la forme B, $x^{2}-4x-32=-32$, soit $x^{2}-4x=0$, puis
$x(x-4)=0$ : $\res{S=\ens{0\,;4}}$.}{le $-32$ s'élimine}
\explique{$E(x)=-36$ : avec la forme A, $(x-2)^{2}-36=-36$, soit $(x-2)^{2}=0$, donc
$x-2=0$ : $\res{S=\ens{2}}$.}{le $-36$ s'élimine}
\rem{La bonne forme est celle où le nombre du second membre s'élimine, ou la forme
factorisée pour une équation $\ldots=0$.}
\end{exo}

\partie{Équations $x^{2}=a$ et équations quotient}

\begin{exo}[Équations $x^{2}=a$]
\explique{On isole le carré, puis : si $a>0$, deux solutions opposées $\sqrt a$ et
$-\sqrt a$ ; si $a=0$, la seule solution $0$ ; si $a<0$, aucune.}{propriété 9 et méthode
13 du cours}
\cas{a)}
\explique{$x^{2}=81$ et $81>0$ : $\res{S=\ens{-9\,;9}}$.}{$9^{2}=81$ et $(-9)^{2}=81$}
\cas{b)}
\explique{$x^{2}=-16$ : un carré n'est jamais négatif, $\res{S=\varnothing}$.}{$-16<0$}
\cas{c)}
\explique{$3x^{2}=48$ donne $x^{2}=16$ : $\res{S=\ens{-4\,;4}}$.}{on divise par $3$ pour
isoler le carré}
\cas{d)}
\explique{$x^{2}-5=0$ donne $x^{2}=5$ : $\res{S=\ens{-\sqrt5\,;\sqrt5}}$.}{$5$ n'est pas un
carré parfait : on garde la racine}
\cas{e)}
\explique{$4x^{2}-25=0$ donne $4x^{2}=25$, puis $x^{2}=\frac{25}{4}$ :
$\res{S=\ens{-\frac52\,;\frac52}}$.}{$\left(\frac52\right)^{2}=\frac{25}{4}$}
\cas{f)}
\explique{$(x+2)^{2}=49$ : le carré porte sur $x+2$, donc $x+2=7$ ou $x+2=-7$, soit $x=5$
ou $x=-9$ : $\res{S=\ens{-9\,;5}}$.}{méthode 13, exemple d)}
\rem{Erreur fréquente : oublier la solution négative ($x^{2}=81$ donne aussi $-9$), ou
écrire $x=\sqrt{-16}$, qui n'existe pas.}
\end{exo}

\begin{exo}[Équations quotient]
\explique{$\dfrac{A}{B}=0$ équivaut à $A=0$ et $B\ne 0$ : on cherche la ou les valeurs
interdites, on résout \og numérateur $=0$\fg{}, puis on retire les valeurs
interdites.}{propriété 9 et méthode 14 du cours}
\cas{a)}
\explique{Valeur interdite : $x+2=0$ pour $x=-2$. Numérateur : $3x-4=0$ donne
$x=\frac43$, qui n'est pas interdit : $\res{S=\ens{\frac43}}$.}{}
\cas{b)}
\explique{Valeur interdite : $3-x=0$ pour $x=3$. Numérateur : $2x+9=0$ donne
$x=-\frac92$ : $\res{S=\ens{-\frac92}}$.}{}
\cas{c)}
\explique{Le dénominateur : $4x^{2}-4=4(x^{2}-1)=4(x+1)(x-1)$ ; il est nul pour $x=-1$ ou
$x=1$, les deux valeurs interdites.}{facteur commun $4$, puis troisième identité}
\begin{calculs}
(2x+3)^{2}-9(x-1)^{2} &= 0 &&\justif{$9(x-1)^{2}=\big[3(x-1)\big]^{2}$}\\
\big[(2x+3)-3(x-1)\big]\big[(2x+3)+3(x-1)\big] &= 0 &&\\
(2x+3-3x+3)(2x+3+3x-3) &= 0 &&\\
(-x+6)\times 5x &= 0 &&
\end{calculs}
\explique{$-x+6=0$ ou $5x=0$, soit $x=6$ ou $x=0$, qui ne sont pas interdits :
$\res{S=\ens{0\,;6}}$.}{}
\cas{d)}
\explique{Le dénominateur : $12-9x=0$ pour $x=\frac{12}{9}=\frac43$, valeur interdite.}{}
\begin{calculs}
9x^{2}-16-2(3x-4) &= 0 &&\justif{$9x^{2}-16=(3x-4)(3x+4)$}\\
(3x-4)(3x+4)-2(3x-4) &= 0 &&\\
(3x-4)\big[(3x+4)-2\big] &= 0 &&\\
(3x-4)(3x+2) &= 0 &&
\end{calculs}
\explique{$3x-4=0$ ou $3x+2=0$, soit $x=\frac43$ ou $x=-\frac23$. Or $\frac43$ est
interdit : $\res{S=\ens{-\frac23}}$.}{on retire la valeur interdite}
\cas{e)}
\explique{Valeur interdite : $0$. Pour $x\ne 0$, on met tout dans le membre de gauche et
au même dénominateur $2x$.}{se ramener à $\frac{A}{B}=0$}
\begin{calculs}
\dfrac{3}{x}+\dfrac12-\dfrac{7}{2x} &= 0 &&\\
\dfrac{3\fois{2}}{x\fois{2}}+\dfrac{1\fois{x}}{2\fois{x}}-\dfrac{7}{2x} &= 0 &&\\
\dfrac{6+x-7}{2x} &= 0 &&\\
\dfrac{x-1}{2x} &= 0 &&
\end{calculs}
\explique{$x-1=0$ donne $x=1$, qui n'est pas interdit : $\res{S=\ens{1}}$.}{vérification :
$3+\frac12=\frac72$}
\rem{Erreur fréquente en d) : garder $\frac43$ dans les solutions. Pour $x=\frac43$, le
quotient n'existe pas.}
\end{exo}

\begin{exo}[Dans des problèmes]
\cas{1)}
\explique{Soit $c$ le côté du carré, en cm ($c>0$) : $c^{2}=30$. Les solutions sont
$\sqrt{30}$ et $-\sqrt{30}$ ; un côté est positif.}{on garde la solution qui a un sens}
\explique{Le côté mesure $\res{\sqrt{30}}$~cm.}{valeur exacte : on ne l'arrondit pas}
\cas{2)}
\explique{Soit $r$ le rayon, en cm : $\pi r^{2}=49\pi$, donc $r^{2}=49$. Comme $r>0$,
$r=7$.}{on divise par $\pi$}
\explique{Le rayon mesure $\res{7}$~cm.}{}
\cas{3)}
\explique{Soit $x$ ce nombre : $x^{2}+7=43$, donc $x^{2}=36$, soit $x=6$ ou $x=-6$.}{}
\explique{Le nombre est $\res{6}$ ou $\res{-6}$.}{vérification : $36+7=43$}
\cas{4)}
\explique{Valeur interdite : $-2$. La fraction est nulle quand $x-6=0$, soit pour
$x=\res{6}$, qui n'est pas interdit.}{équation quotient}
\explique{Elle vaut $1$ si $x-6=x+2$ (produit en croix), soit $0x=8$ : c'est impossible.
La fraction \res{ne vaut jamais $1$}.}{le numérateur est toujours plus petit de $8$ que le
dénominateur}
\end{exo}

\partie{Mise en équation}

\begin{exo}[Quatre problèmes courts]
\explique{Pour chaque problème : on annonce l'inconnue (\og Soit $x$ \ldots\fg), on traduit
l'énoncé par une équation, on la résout, on vérifie que le résultat a un sens, puis on
conclut par une phrase.}{méthode 15 du cours}
\cas{1)}
\explique{Soit $x$ l'âge de Zoé, en années. Elle ajoute $9$ : $x+9$ ; elle multiplie le
résultat par $3$ : $3(x+9)$.}{le \og résultat\fg{} entier est multiplié : parenthèses}
\begin{calculs}
3(x+9) &= 75 &&\\
3x+27 &= 75 &&\\
3x &= 48 &&\\
x &= 16 &&
\end{calculs}
\explique{Zoé a $\res{16}$ ans.}{vérification : $16+9=25$ et $25\times 3=75$}
\cas{2)}
\explique{Soit $x$ le nombre choisi par Lina : $\times 3$ donne $3x$, $-4$ donne $3x-4$,
puis $\times 5$ donne $5(3x-4)$.}{on suit le programme étape par étape}
\begin{calculs}
5(3x-4) &= 205 &&\\
15x-20 &= 205 &&\\
15x &= 225 &&\\
x &= 15 &&
\end{calculs}
\explique{Lina avait choisi $\res{15}$.}{vérification : $45-4=41$ et $41\times 5=205$}
\cas{3)}
\explique{Soit $x$ la contenance de la cuve, en litres. Remplie au quart, elle contient
$\frac{x}{4}$ ; aux deux tiers, $\frac{2x}{3}$.}{\og le quart de $x$\fg{} : $\frac14\times x$}
\begin{calculs}
\dfrac{x}{4}+15 &= \dfrac{2x}{3} &&\justif{on ajoute $15$~L}\\
3x+180 &= 8x &&\justif{chaque terme $\times 12$}\\
180 &= 5x &&\\
x &= 36 &&
\end{calculs}
\explique{La cuve contient $\res{36}$~L.}{vérification : $9+15=24$, et $24$ est les deux
tiers de $36$}
\cas{4)}
\explique{Soit $n$ le plus petit des deux entiers impairs : le suivant est $n+2$.}{deux
impairs consécutifs diffèrent de $2$, pas de $1$}
\begin{calculs}
n+(n+2) &= 152 &&\\
2n+2 &= 152 &&\\
2n &= 150 &&\\
n &= 75 &&
\end{calculs}
\explique{Les deux entiers sont $\res{75}$ et $\res{77}$.}{$75$ est bien impair ;
vérification : $75+77=152$}
\end{exo}

\begin{exo}[Entiers et âges]
\cas{1)}
\explique{Soit $n$ le premier entier : les deux suivants sont $n+1$ et $n+2$. Le
troisième est $n+2$, son triple $3(n+2)$.}{on nomme le plus petit}
\begin{calculs}
n+3(n+2) &= 98 &&\justif{le premier, augmenté du triple du troisième}\\
n+3n+6 &= 98 &&\\
4n &= 92 &&\\
n &= 23 &&
\end{calculs}
\explique{Les trois entiers sont $\res{23\,;24}$ et $\res{25}$.}{vérification :
$23+3\times 25=98$}
\cas{2)}
\explique{Soit $x$ le nombre d'années cherché. Dans $x$ ans, les sœurs auront $30+x$,
$18+x$ et $4+x$ ans.}{chacune vieillit de $x$ ans}
\begin{calculs}
30+x &= (18+x)+(4+x) &&\justif{l'aînée, la somme des deux autres}\\
30+x &= 22+2x &&\\
30-22 &= 2x-x &&\\
x &= 8 &&
\end{calculs}
\explique{Ce sera dans $\res{8}$ ans.}{vérification : $38=26+12$}
\rem{Erreur fréquente : écrire $30+x=18+4+x$, en ne faisant vieillir qu'une des deux
cadettes.}
\end{exo}

\begin{exo}[Stocks, jardin et concert]
\cas{1)}
\explique{Soit $b$ le nombre de paires de bottes : il y a $3b$ paires de baskets et $b+5$
paires de sandales.}{on part de celle qui sert à exprimer les autres}
\begin{calculs}
b+3b+(b+5) &= 95 &&\\
5b+5 &= 95 &&\\
5b &= 90 &&\\
b &= 18 &&
\end{calculs}
\explique{Le magasin reçoit $\res{18}$ paires de bottes, $\res{54}$ paires de baskets et
$\res{23}$ paires de sandales.}{vérification : $18+54+23=95$}
\cas{2)}
\explique{Soit $x$ l'aire du jardin, en m$^{2}$. Le verger est ce qui reste quand on
enlève le quart et le tiers : $x-\frac{x}{4}-\frac{x}{3}$.}{}
\begin{calculs}
x-\dfrac{x}{4}-\dfrac{x}{3} &= 75 &&\\
12x-3x-4x &= 900 &&\justif{chaque terme $\times 12$}\\
5x &= 900 &&\\
x &= 180 &&
\end{calculs}
\explique{L'aire du jardin est de $\res{180}$~m$^{2}$.}{vérification : $180-45-60=75$}
\cas{3)}
\explique{Soit $x$ le nombre de places assises : il y a $800-x$ places debout.}{les deux
nombres font $800$ en tout}
\begin{calculs}
40x+25(800-x) &= 26\,000 &&\justif{la recette}\\
40x+20\,000-25x &= 26\,000 &&\\
15x &= 6\,000 &&\\
x &= 400 &&
\end{calculs}
\explique{$\res{400}$ places assises ont été vendues.}{vérification :
$16\,000+10\,000=26\,000$}
\end{exo}

\partie{Problèmes de géométrie}

\begin{exo}[Exprimer des longueurs et des aires]
\cas{1)}
\explique{$GO$ est la somme des deux morceaux : $GO=\res{x+7}$.}{on lit les cotes}
\explique{$[RT]$, de longueur $x$, est partagé en cinq parts égales de $\frac{x}{5}$, et
$[RS]$ en contient deux : $RS=\res{\dfrac{2x}{5}}$.}{$2\times\frac{x}{5}=\frac{2x}{5}$}
\cas{2)}
\explique{La partie colorée est le grand rectangle de dimensions $x+6$ et $6$, moins le
rectangle blanc de dimensions $6$ et $x$.}{on découpe la figure}
\begin{calculs}
\mathcal{A} &= 6(x+6)-6x &&\\
\mathcal{A} &= 6x+36-6x &&\\
\mathcal{A} &= \res{36} &&
\end{calculs}
\explique{L'aire de la partie colorée vaut $36$ quelle que soit la valeur de $x$ : elle ne
dépend pas de $x$.}{c'est l'aire d'un carré de côté $6$, déplacé en partie}
\end{exo}

\begin{exo}[Un trapèze et un triangle]
\cas{1)}
\explique{L'aire d'un trapèze est $\dfrac{(B+b)\times h}{2}$, avec ici $B=3x$ ; $b=x$ et
$h=x+3$.}{formule du collège}
\begin{calculs}
\mathcal{A}_{1} &= \dfrac{(3x+x)(x+3)}{2} &&\\
\mathcal{A}_{1} &= \dfrac{4x(x+3)}{2} &&\\
\mathcal{A}_{1} &= 2x(x+3) &&\justif{$4\div 2=2$}\\
\mathcal{A}_{1} &= \res{2x^{2}+6x} &&
\end{calculs}
\explique{L'aire d'un triangle est $\dfrac{\text{base}\times\text{hauteur}}{2}$ : la base
mesure $6+x$, la hauteur $2x$.}{}
\begin{calculs}
\mathcal{A}_{2} &= \dfrac{(6+x)\times 2x}{2} &&\\
\mathcal{A}_{2} &= (6+x)\times x &&\justif{$2x\div 2=x$}\\
\mathcal{A}_{2} &= \res{x^{2}+6x} &&
\end{calculs}
\cas{2)}
\begin{calculs}
\mathcal{A}_{1}+\mathcal{A}_{2} &= 2x^{2}+6x+x^{2}+6x &&\\
\mathcal{A}_{1}+\mathcal{A}_{2} &= 3x^{2}+12x &&\\
\mathcal{A}_{1}+\mathcal{A}_{2} &= 3x(x+4) &&\justif{facteur commun $3x$}
\end{calculs}
\explique{La somme des aires est l'aire d'un rectangle de côtés $3x$ et
$\res{x+4}$.}{factoriser par $3x$ fait apparaître l'autre côté}
\end{exo}

\begin{exo}[Deux solides]
\explique{Volume du pavé droit : longueur $\times$ largeur $\times$ hauteur.}{}
\begin{calculs}
V_{1} &= (2x+6)\times 2x\times x &&\\
V_{1} &= (2x+6)\times 2x^{2} &&\\
V_{1} &= 4x^{3}+12x^{2} &&
\end{calculs}
\explique{Volume du prisme droit : aire de la base $\times$ longueur. La base est un
triangle rectangle de côtés de l'angle droit $2x$ et $4x$.}{}
\begin{calculs}
\mathcal{B} &= \dfrac{2x\times 4x}{2} &&\\
\mathcal{B} &= 4x^{2} &&\\
V_{2} &= 4x^{2}(x+3) &&\justif{la longueur du prisme est $x+3$}\\
V_{2} &= 4x^{3}+12x^{2} &&
\end{calculs}
\explique{Pour tout $x>0$, $V_{1}$ et $V_{2}$ valent $4x^{3}+12x^{2}$ : les deux solides ont
\res{le même volume}.}{on développe les deux volumes pour les comparer}
\end{exo}

\begin{exo}[Une aire donnée]
\explique{$ABCD$ est un rectangle de longueur $AB=6+x$ et de largeur $BC=4$ (en cm, avec
$x>0$).}{on exprime les longueurs}
\begin{calculs}
4(6+x) &= 30 &&\justif{l'aire de $ABCD$}\\
24+4x &= 30 &&\\
4x &= 6 &&\\
x &= 1{,}5 &&\justif{$1{,}5>0$ : possible}
\end{calculs}
\explique{Il faut $x=\res{1{,}5}$~cm.}{vérification : $4\times 7{,}5=30$}
\end{exo}

\begin{exo}[Un carré dans un carré]
\explique{La bande a la largeur $x$ de chaque côté : le grand carré a pour côté
$x+3+x=3+2x$ (en cm, avec $x>0$).}{on place $x$ des deux côtés du petit carré}
\cas{1)}
\begin{calculs}
4(3+2x) &= 36 &&\justif{périmètre du grand carré}\\
12+8x &= 36 &&\\
8x &= 24 &&\\
x &= 3 &&
\end{calculs}
\explique{Il faut $x=\res{3}$~cm.}{vérification : côté $9$, périmètre $36$}
\cas{2)}
\explique{$(3+2x)^{2}=49$ : le carré porte sur $3+2x$, donc $3+2x=7$ ou $3+2x=-7$, soit
$x=2$ ou $x=-5$.}{méthode 13, exemple d)}
\explique{Une largeur est positive : il faut $x=\res{2}$~cm.}{on rejette $-5$}
\cas{3)}
\explique{Aire de la bande : grand carré moins petit carré, $(3+2x)^{2}-9$. Le triple de
l'aire du petit carré vaut $3\times 9=27$.}{}
\begin{calculs}
(3+2x)^{2}-9 &= 27 &&\\
(3+2x)^{2} &= 36 &&
\end{calculs}
\explique{$3+2x=6$ ou $3+2x=-6$, soit $x=1{,}5$ ou $x=-4{,}5$. Il faut $x=\res{1{,}5}$~cm.}{vérification :
côté $6$, bande $36-9=27$}
\end{exo}

\begin{exo}[Un bassin dans une parcelle]
\cas{1)}
\explique{Le bassin est un carré de côté $x-2$ : la bordure est la parcelle moins le
bassin, d'aire $x^{2}-(x-2)^{2}$. L'expression \res{d)} convient.}{on découpe la figure}
\begin{calculs}
x^{2}-(x-2)^{2} &= x^{2}-\big[x^{2}-4x+4\big] &&\\
x^{2}-(x-2)^{2} &= 4x-4 &&
\end{calculs}
\explique{L'expression \res{c)} convient donc aussi. Les autres ne conviennent pas :
a) vaut $4x+4$ ; b) vaut $4x$ ; e) vaut $x^{2}-2x$.}{pour $x=3$ : la bordure vaut
$9-1=8$ ; a) donne $16$, b) $12$, e) $3$}
\cas{2)}
\explique{$(x-2)^{2}=36$ donne $x-2=6$ ou $x-2=-6$, soit $x=8$ ou $x=-4$.}{méthode 13}
\explique{Comme $x>2$, la parcelle a pour côté $x=\res{8}$~m.}{vérification : bassin de
côté $6$, d'aire $36$}
\end{exo}

\partie{Isoler une grandeur}

\begin{exo}[Échauffement]
\explique{On traite $A$ comme l'inconnue d'une équation, les autres lettres comme des
nombres connus.}{définition 5 et méthode 17 du cours}
\cas{a)}
\begin{calculs}
A+3 &= B &&\\
A &= \res{B-3} &&\justif{on soustrait $3$}
\end{calculs}
\cas{b)}
\begin{calculs}
5-A &= B &&\\
5 &= B+A &&\justif{on ajoute $A$ pour le rendre positif}\\
A &= \res{5-B} &&\justif{on soustrait $B$}
\end{calculs}
\cas{c)}
\begin{calculs}
4A &= B &&\\
A &= \res{\dfrac{B}{4}} &&\justif{on divise par $4$}
\end{calculs}
\cas{d)}
\begin{calculs}
\dfrac{A}{3} &= B &&\\
A &= \res{3B} &&\justif{on multiplie par $3$}
\end{calculs}
\cas{e)}
\begin{calculs}
\dfrac{6}{A} &= B &&\justif{produit en croix avec $\frac{B}{1}$}\\
6 &= AB &&\\
A &= \res{\dfrac{6}{B}} &&\justif{on divise par $B$, non nul}
\end{calculs}
\cas{f)}
\begin{calculs}
2A+B &= C &&\\
2A &= C-B &&\justif{on soustrait $B$}\\
A &= \res{\dfrac{C-B}{2}} &&\justif{on divise par $2$ tout le membre de droite}
\end{calculs}
\rem{Erreur fréquente en f) : $A=C-\frac{B}{2}$. On divise \emph{tout} $C-B$ par $2$.}
\end{exo}

\begin{exo}[Formules de géométrie et de physique]
\cas{1)}
\begin{calculs}
V &= L\,\ell\,h &&\justif{on recopie}\\
\ell &= \res{\dfrac{V}{L\,h}} &&\justif{on divise par $L\,h$}
\end{calculs}
\cas{2)}
\begin{calculs}
\mathcal{A} &= \dfrac{b\times h}{2} &&\\
2\mathcal{A} &= b\times h &&\justif{on multiplie par $2$}\\
h &= \res{\dfrac{2\mathcal{A}}{b}} &&\justif{on divise par $b$}
\end{calculs}
\cas{3)}
\begin{calculs}
U &= E-rI &&\\
rI &= E-U &&\justif{on ajoute $rI$ et on soustrait $U$}\\
r &= \res{\dfrac{E-U}{I}} &&\justif{on divise par $I$}
\end{calculs}
\cas{4)}
\begin{calculs}
L\,H-\dfrac{V}{d} &= 0 &&\\
L\,H &= \dfrac{V}{d} &&\justif{on ajoute $\frac{V}{d}$}\\
L\,H\times d &= V &&\justif{produit en croix avec $\frac{L\,H}{1}$}\\
d &= \res{\dfrac{V}{L\,H}} &&\justif{on divise par $L\,H$}
\end{calculs}
\cas{5)}
\begin{calculs}
C &= C_{0}(1+t) &&\\
C_{0} &= \res{\dfrac{C}{1+t}} &&\justif{on divise par $1+t$ ; on ne développe pas}
\end{calculs}
\cas{6)}
\begin{calculs}
\mathcal{A} &= \pi r^{2} &&\\
r^{2} &= \dfrac{\mathcal{A}}{\pi} &&\justif{on divise par $\pi$}\\
r &= \res{\sqrt{\dfrac{\mathcal{A}}{\pi}}} &&\justif{$r$ est positif : on garde la racine positive}
\end{calculs}
\end{exo}

\begin{exo}[Plus difficile]
\cas{1)}
\begin{calculs}
\dfrac1f &= \dfrac1a+\dfrac1b &&\justif{on recopie}\\
\dfrac1f &= \dfrac{b}{ab}+\dfrac{a}{ab} &&\justif{dénominateur commun $ab$}\\
\dfrac1f &= \dfrac{a+b}{ab} &&\\
f\times(a+b) &= ab &&\justif{produit en croix}\\
f &= \res{\dfrac{ab}{a+b}} &&\justif{on divise par $a+b$}
\end{calculs}
\rem{Erreur fréquente : $f=a+b$. L'inverse d'une somme n'est pas la somme des inverses.}
\cas{2)}
\begin{calculs}
F &= 6{,}67\times 10^{-11}\,\dfrac{m\,m'}{d^{2}} &&\\
F\times d^{2} &= 6{,}67\times 10^{-11}\,m\,m' &&\justif{on multiplie par $d^{2}$}\\
d^{2} &= \dfrac{6{,}67\times 10^{-11}\,m\,m'}{F} &&\justif{on divise par $F$}\\
d &= \res{\sqrt{\dfrac{6{,}67\times 10^{-11}\,m\,m'}{F}}} &&\justif{$d$ est positif}
\end{calculs}
\end{exo}

\begin{exo}[Isoler, puis calculer]
\cas{1)}
\begin{calculs}
4x+2y-3z &= 10 &&\\
2y &= 10-4x+3z &&\justif{on soustrait $4x$ et on ajoute $3z$}\\
y &= \res{\dfrac{10-4x+3z}{2}} &&\justif{on divise par $2$}
\end{calculs}
\begin{calculs}
y &= \dfrac{10-4\times 1+3\times 2}{2} &&\justif{pour $x=1$ et $z=2$}\\
y &= \dfrac{12}{2} &&\\
y &= \res{6} &&\justif{vérification : $4+12-6=10$}
\end{calculs}
\cas{2)}
\explique{Avec l'exercice 52 (question 1), $\ell=\dfrac{V}{L\,h}$, avec $V=72$, $L=6$
et $h=3$.}{on isole d'abord, puis on substitue}
\begin{calculs}
\ell &= \dfrac{72}{6\times 3} &&\\
\ell &= \dfrac{72}{18} &&\\
\ell &= 4 &&
\end{calculs}
\explique{La largeur du pavé est de $\res{4}$~cm.}{vérification : $6\times 4\times 3=72$}
\cas{3)}
\explique{Avec l'exercice 52 (question 6), $r=\sqrt{\dfrac{\mathcal{A}}{\pi}}$, avec
$\mathcal{A}=80$.}{}
\begin{calculs}
r &= \sqrt{\dfrac{80}{\pi}} &&\\
r &\approx 5{,}046 &&\justif{à la calculatrice}
\end{calculs}
\explique{Le rayon mesure environ $\res{5{,}0}$~cm (au millimètre près).}{le chiffre des centièmes est $4$ :
on arrondit à $5{,}0$}
\end{exo}

\partiesansnum{Vers le devoir commun}

\begin{exo}[Maîtrise des calculs]
\cas{1) $A(x)$}
\begin{calculs}
A(x) &= 3(2x-5)+(3x+4)(2x-1) &&\justif{une distributivité, puis une double distributivité}\\
A(x) &= 6x-15+6x^{2}-3x+8x-4 &&\\
A(x) &= \res{6x^{2}+11x-19} &&
\end{calculs}
\cas{1) $B(x)$}
\begin{calculs}
B(x) &= (x+2)(-3x+4)-(5-x) &&\\
B(x) &= -3x^{2}+4x-6x+8-5+x &&\justif{le $-$ change les signes de $5-x$}\\
B(x) &= \res{-3x^{2}-x+3} &&
\end{calculs}
\cas{2) $C(x)$}
\begin{calculs}
C(x) &= 8x^{2}-5x &&\justif{$x$ est commun}\\
C(x) &= \res{x(8x-5)} &&
\end{calculs}
\cas{2) $D(x)$}
\explique{$9x^{2}=(3x)^{2}$ et $25=5^{2}$ ; $2\times 3x\times 5=30x$ est le terme restant.
Donc $D(x)=\res{(3x+5)^{2}}$.}{carré parfait}
\cas{2) $E(x)$}
\begin{calculs}
E(x) &= (2x-5)-3(2x-5)^{2} &&\justif{le « 1 » caché : $(2x-5)=1\times(2x-5)$}\\
E(x) &= (2x-5)\big[1-3(2x-5)\big] &&\\
E(x) &= (2x-5)(1-6x+15) &&\\
E(x) &= \res{(2x-5)(16-6x)} &&
\end{calculs}
\rem{Et $16-6x=2(8-3x)$ : $E(x)=2(2x-5)(8-3x)$.}
\cas{2) $F(x)$}
\begin{calculs}
F(x) &= (x+2)(-2x+3)-(x+2)^{2} &&\justif{facteur commun $(x+2)$}\\
F(x) &= (x+2)\big[(-2x+3)-(x+2)\big] &&\\
F(x) &= (x+2)(-2x+3-x-2) &&\\
F(x) &= \res{(x+2)(-3x+1)} &&
\end{calculs}
\end{exo}

\begin{exo}[Une fonction, trois écritures]
\cas{1)}
\begin{calculs}
f(x) &= (x+3)^{2}-16 &&\justif{forme de l'énoncé}\\
f(x) &= x^{2}+6x+9-16 &&\justif{$(a+b)^{2}$ avec $a=x$ et $b=3$}\\
f(x) &= \res{x^{2}+6x-7} &&\justif{forme développée}
\end{calculs}
\begin{calculs}
f(x) &= (x+3)^{2}-4^{2} &&\justif{$a^{2}-b^{2}$ avec $a=x+3$ et $b=4$}\\
f(x) &= \big[(x+3)-4\big]\big[(x+3)+4\big] &&\\
f(x) &= \res{(x-1)(x+7)} &&\justif{forme factorisée}
\end{calculs}
\cas{2)}
\explique{$f(0)=0+0-7=\res{-7}$.}{forme développée}
\begin{calculs}
f\big(\tfrac12\big) &= \big(\tfrac12+3\big)^{2}-16 &&\justif{forme de l'énoncé : $\frac12+3=\frac72$}\\
f\big(\tfrac12\big) &= \tfrac{49}{4}-\tfrac{64}{4} &&\justif{$16=\frac{64}{4}$}\\
f\big(\tfrac12\big) &= \res{-\tfrac{15}{4}} &&
\end{calculs}
\explique{Antécédents de $0$ : $(x-1)(x+7)=0$, soit $x=1$ ou $x=-7$. Ce sont $\res{1}$ et
$\res{-7}$.}{forme factorisée}
\explique{Antécédents de $-7$ : $x^{2}+6x-7=-7$, soit $x^{2}+6x=0$, puis $x(x+6)=0$. Ce
sont $\res{0}$ et $\res{-6}$.}{forme développée : le $-7$ s'élimine}
\explique{Antécédent de $-16$ : $(x+3)^{2}-16=-16$, soit $(x+3)^{2}=0$, donc $x=-3$. C'est
$\res{-3}$.}{forme de l'énoncé : le $-16$ s'élimine}
\end{exo}

\begin{exo}[Lire, puis démontrer]
\cas{1)}
\explique{La courbe coupe l'axe des abscisses aux points d'abscisses $-1$ et $4$ : on
conjecture que $f(x)=0$ a pour solutions $\res{-1}$ et $\res{4}$.}{points de la courbe
d'ordonnée $0$}
\explique{La courbe passe à l'ordonnée $-4$ aux abscisses $0$ et $3$ : on conjecture
que $f(x)=-4$ a pour solutions $\res{0}$ et $\res{3}$.}{on trace la droite horizontale
d'ordonnée $-4$}
\cas{2)}
\begin{calculs}
f(x) &= (2x-3)(x+1)-(x+1)^{2} &&\\
f(x) &= 2x^{2}+2x-3x-3-\big[x^{2}+2x+1\big] &&\justif{le carré entre crochets}\\
f(x) &= 2x^{2}-x-3-x^{2}-2x-1 &&\\
f(x) &= \res{x^{2}-3x-4} &&
\end{calculs}
\begin{calculs}
f(x) &= (2x-3)(x+1)-(x+1)^{2} &&\justif{on repart de l'énoncé : facteur commun $(x+1)$}\\
f(x) &= (x+1)\big[(2x-3)-(x+1)\big] &&\\
f(x) &= (x+1)(2x-3-x-1) &&\\
f(x) &= \res{(x+1)(x-4)} &&
\end{calculs}
\cas{3)}
\explique{$f(2)=4-6-4=\res{-6}$.}{forme développée}
\explique{$f(x)=0$ : $(x+1)(x-4)=0$, soit $x=-1$ ou $x=4$, toutes deux dans $\intff{-2}{5}$ :
$\res{S=\ens{-1\,;4}}$.}{forme factorisée}
\explique{$f(x)=-4$ : $x^{2}-3x-4=-4$, soit $x^{2}-3x=0$, puis $x(x-3)=0$ :
$\res{S=\ens{0\,;3}}$.}{forme développée : le $-4$ s'élimine}
\explique{Les solutions trouvées par le calcul sont celles lues sur le graphique : les deux
conjectures sont démontrées.}{une lecture graphique se démontre par le calcul}
\end{exo}

\begin{exo}[Équations]
\cas{a)}
\begin{calculs}
\dfrac23x-3 &= 5 &&\\
\dfrac23x &= 8 &&\justif{on ajoute $3$}\\
x &= 8\times\dfrac32 &&\justif{on multiplie par l'inverse de $\frac23$}\\
x &= 12 &&
\end{calculs}
\explique{$\res{S=\ens{12}}$.}{vérification : $8-3=5$}
\cas{b)}
\begin{calculs}
\dfrac74-\dfrac{x}{3} &= 1 &&\justif{chaque terme $\times 12$}\\
21-4x &= 12 &&\\
9 &= 4x &&\\
x &= \dfrac94 &&
\end{calculs}
\explique{$\res{S=\ens{\frac94}}$.}{vérification : $\frac74-\frac34=1$}
\cas{c)}
\begin{calculs}
4x+3 &= 1-3(2-x) &&\\
4x+3 &= 1-6+3x &&\justif{$-3\times(-x)=+3x$}\\
4x+3 &= 3x-5 &&\\
x &= -8 &&
\end{calculs}
\explique{$\res{S=\ens{-8}}$.}{vérification : $-29=-29$}
\cas{d)}
\begin{calculs}
(3x-4)^{2}-(x+6)^{2} &= 0 &&\justif{$a=3x-4$ et $b=x+6$}\\
\big[(3x-4)-(x+6)\big]\big[(3x-4)+(x+6)\big] &= 0 &&\\
(2x-10)(4x+2) &= 0 &&
\end{calculs}
\explique{$2x-10=0$ ou $4x+2=0$ : $\res{S=\ens{-\frac12\,;5}}$.}{}
\cas{e)}
\explique{$9x^{2}-12x+4=(3x-2)^{2}$ ($a=3x$ et $b=2$ ; $2ab=12x$) et
$4(x+2)^{2}=\big[2(x+2)\big]^{2}$.}{on reconnaît deux carrés}
\begin{calculs}
9x^{2}-12x+4 &= 4(x+2)^{2} &&\\
(3x-2)^{2}-\big[2(x+2)\big]^{2} &= 0 &&\justif{tout à gauche}\\
\big[(3x-2)-2(x+2)\big]\big[(3x-2)+2(x+2)\big] &= 0 &&\\
(x-6)(5x+2) &= 0 &&
\end{calculs}
\explique{$x-6=0$ ou $5x+2=0$ : $\res{S=\ens{-\frac25\,;6}}$.}{}
\end{exo}

\begin{exo}[Points d'intersection de deux courbes]
\explique{Un point commun aux deux courbes a la même abscisse $x$ et la même ordonnée : on
résout $f(x)=g(x)$, puis on calcule l'ordonnée $g(x)$.}{chapitre 2}
\cas{a)}
\begin{calculs}
x^{2}+2x-3 &= 2x+6 &&\\
x^{2}-9 &= 0 &&\justif{les $2x$ s'éliminent}\\
(x-3)(x+3) &= 0 &&
\end{calculs}
\explique{$x=3$ ou $x=-3$ ; $g(3)=12$ et $g(-3)=0$. Les points d'intersection sont
$\res{(3\,;12)}$ et $\res{(-3\,;0)}$.}{contrôle : $f(3)=9+6-3=12$}
\cas{b)}
\begin{calculs}
3x^{2}-x+1 &= 5x+1 &&\\
3x^{2}-6x &= 0 &&\\
3x(x-2) &= 0 &&
\end{calculs}
\explique{$x=0$ ou $x=2$ ; $g(0)=1$ et $g(2)=11$. Les points d'intersection sont
$\res{(0\,;1)}$ et $\res{(2\,;11)}$.}{contrôle : $f(2)=12-2+1=11$}
\end{exo}

\begin{exo}[Deux problèmes]
\cas{1)}
\explique{Soit $x$ le côté du carré de départ, en cm ($x>0$) : son aire est $x^{2}$ ; le
nouveau carré a pour aire $(x+3)^{2}$.}{}
\begin{calculs}
(x+3)^{2} &= x^{2}+51 &&\\
x^{2}+6x+9 &= x^{2}+51 &&\\
6x &= 42 &&\justif{les $x^{2}$ s'éliminent}\\
x &= 7 &&
\end{calculs}
\explique{L'aire du carré de départ est $7^{2}=\res{49}$~cm$^{2}$.}{vérification :
$100=49+51$}
\cas{2)}
\explique{$M$ est sur $[AB]$ : $0<x\le 6$. Le carré a pour aire $x^{2}$. Le rectangle a
pour côtés $NQ=6-x$ et $QC=9-x$, donc pour aire $(6-x)(9-x)=54-15x+x^{2}$.}{comme dans
la méthode 16}
\begin{calculs}
x^{2} &= 54-15x+x^{2} &&\\
15x &= 54 &&\justif{les $x^{2}$ s'éliminent}\\
x &= 3{,}6 &&\justif{$0<3{,}6\le 6$ : possible}
\end{calculs}
\explique{Il faut placer $M$ à $\res{3{,}6}$~cm de $A$.}{vérification : $3{,}6^{2}=12{,}96$
et $2{,}4\times 5{,}4=12{,}96$}
\end{exo}

\end{document}
